% vim: tw=80 noai
\documentclass[normaltoc,capchap,capsec,times]{abnt}
\usepackage[utf8x]{inputenc}
\usepackage[T1]{fontenc}
\usepackage[brazil]{babel}
\usepackage[alf]{abntcite}
\usepackage[ordem=alf]{tabela-simbolos}
\usepackage{url}
\usepackage[pdftex]{graphicx}
\usepackage{listings}
\usepackage{verbatim}
\usepackage{subfigure}
\usepackage{multicol}
\usepackage{multirow}
\usepackage{framed}
%\usepackage{algorithm2e}
\usepackage{amssymb, amsmath}


\def\lstlistingname{Listagem}

%%%%%%%%%%%%%%%%%%%%%%%%%%%%%%%%%%%%%%%%%%%%%%%%%%%%
% Dados da monografia
%%%%%%%%%%%%%%%%%%%%%%%%%%%%%%%%%%%%%%%%%%%%%%%%%%%%

\newcommand{\meunome}{Aline Duarte Bessa}
\newcommand{\meutitulo}{Classificação de documentos de acordo com pontos de vista sobre um tema}
\newcommand{\meusubtitulo}{Aspectos teóricos e práticos}
\newcommand{\meuano}{2010}
\newcommand{\meuorientador}{Orientador: Alexandre Tachard Passos}
\newcommand{\meucoorientador}{Co-orientador: Luciano Porto Barreto}

%%%%%%%%%%%%%%%%%%%%%%%%%%%%%%%%%%%%%%%%%%%%%%%%%%%%

%% O comando \obs aqui definido permite que o autor faca anotacoes na
%% monografia que aparecem no PDF gerado. Para ativar o comando, descomente
%% a primeira linha e comente a segunda.
%% Exemplo de uso: \obs{Preciso melhorar este parágrafo...}

%\newcommand{\obs}[1]{\underline{\textbf{OBSERVAÇÃO}}: \emph{#1}}
\newcommand{\obs}[1]{}

\def\ordfem{\mbox{\raise .35em \hbox{\underline{\scriptsize a}\ }}}
\def\ordmasc{\mbox{\raise .35em \hbox{\underline{\scriptsize o}\ }}}
\def\profa{Prof\ordfem.}

%%%%%%%%%%%%%%%%%%%%%%%%%%%%%%%%%%%%%%%%%%%%%%%%%%%%

\begin{document}

% Capa com Brasão

\begin{titlepage}
   \begin{center}
      %logotipo
               \includegraphics{brasaoUFBA} \\
  %\vspace{0.7in}
              \centering{ 
        \bf{
        \LARGE{
    \uppercase{UNIVERSIDADE FEDERAL DA BAHIA} \\
        }
        \Large {
                    \uppercase{INSTITUTO DE MATEMÁTICA} \\
        }
                   \large {
                       \uppercase{DEPARTAMENTO DE CIÊNCIA DA COMPUTAÇÃO} \\
                  }
              } }
   \end{center}
   \vfill
   \begin{center}
       \bf{
       \large{\uppercase{\meunome}  \\  }
       }
   \end{center}
   \vspace{0.2in}
   \begin{center}
       \bf{
         \LARGE{ \uppercase{\meutitulo} } \\
         \Large{ \uppercase{\meusubtitulo} }
         \obs{\\ \Large{Esta versão da monografia contém comentários do autor.
          Para removê-los, redefina o comando LaTeX \texttt{obs}.}}
       }
   \end{center}

   \vfill
   \hspace{\stretch{1}}
   \vfill
   \begin{center}
      \normalsize{
          Salvador \\
          \meuano
       }
   \end{center}

\end{titlepage}

%comando abaixo cria uma capa redundante, mas como a capa com brasão foi 
% feita 'manualmente', não faz sentido usar este comando:
%\capa



%\folhaderosto
% o comando acima foi comentado para não criar uma folha de 
% rosto redundante, já que ela feita 'manualmente' abaixo

\begin{titlepage}
 \vfill
 \begin{center}
   {\large \uppercase{ \bf{ \meunome\ } } } \\[7cm]
   {\Huge \uppercase{ \bf{ \meutitulo\ } } }\\[1cm]
   \vfill
   \hspace{.45\textwidth} % posicionando a minipage
   \begin{minipage}{.5\textwidth}
     \begin{espacosimples}
       \bf{
  Monografia apresentada ao Curso de graduação em Ciência da Computação, 
  Departamento de Ciência da Computação, Instituto de Matemática,  Universidade Federal da 
  Bahia, como requisito parcial para obtenção do grau de Bacharel em  Ciência da Computação. \\ 
       }      
     \end{espacosimples}
     \begin{espacosimples}    
       \meuorientador \\
       \meucoorientador
     \end{espacosimples}
   \end{minipage}
   \vfill
   Salvador \\
   \meuano
 \end{center}
\end{titlepage}

\chapter*{Agradecimentos}

À UFBA, por viabilizar esta monografia.

A meu orientador e amigo Alexandre Passos, por me ajudar a definir o tema, me motivar e participar ativamente de todo o processo.

A meu co-orientador Luciano Porto, pelas sugestões e por acompanhar esta monografia do começo ao fim.

À minha mãe, por ser meu baluarte durante todos esses anos. A meu pai, por me fazer ver tudo de forma mais simples.

A Andrea e Lucas, jornalistas que me ajudaram a selecionar bons \emph{sites} de política brasileira, fundamentais para essa monografia. E a Marcelo, por gostar tanto de política e analisar esse trabalho sob a ótica de alguém não-envolvido.

A todos os meus amigos, pelos momentos deliciosos nos fins de semana e compreensão com meus sumiços.



\begin{resumo}
A classificação de documentos de acordo com seus pontos de vista sobre um tema é um tópico de pesquisa que vem recebendo atenção crescente. Dado um conjunto de documentos sobre a política dos Estados Unidos, por exemplo, esse tipo de classificação pode ser empregado 
para identificar quais deles defendem o ponto de vista \emph{liberal} e quais defendem o ponto de vista 
\emph{conservador}. Esse tipo de classificação colabora com a análise de documentos argumentativos, especialmente quando não é possível lê-los um por  um. Devido ao caráter desses documentos, esse tipo de classificação pode fazer parte de estudos interdisciplinares, envolvendo áreas como Ciência Política e Comunicação. 

Diante disso, essa monografia apresenta uma revisão de trabalhos que tratam desse tópico, indicando quais são os principais classificadores utilizados, quais aspectos dos documentos podem ser explorados na identificação de seus pontos de vista, dentre outros aspectos associados a esse tópico de pesquisa. Em seguida, é feito um estudo de caso envolvendo documentos sobre a política brasileira escritos em 2010, ano de eleição presidencial. Eles são classificados de acordo com os pontos de vista \emph{pró-oposição} e \emph{pró-situação}, e as palavras mais enfatizadas por cada ponto de vista também são discutidas. 


%    Diante disso, estudos com o intuito de extrair opiniões e pontos de vista da Web, e interpretá-
%los automaticamente, tornaram-se mais frequentes na área de Ciência da Computação. Juntos,

%, que visa a separar documentos de acordo com seus diferentes posicionamentos sobre um tema


%e outros aspectos associados a 

\textbf{Palavras-chave:}
Mineração de opinião,
classificação, 
política brasileira,
ponto de vista.
\end{resumo}

% O abstract e' opcional.
\begin{abstract}
Documents classification according to their viewpoints is a research topic that has been receiving a growing attention. Given a set of documents about the United States politics, for example, this kind of classification may be useful to identify those that defend a \emph{liberal} viewpoint and those that defend a \emph{conservative} one. This kind of classification is helpful in analyzing argumentative documents, specially when it is not possible to read them one by one. Due to the content of these documents, this kind of classification may be part of interdisciplinary studies, involving areas such as Political Sciences or Communication.

In the view of this, this monograph presents a review of papers that discuss this topic, indicating which classifiers are the most used, which aspects of the documents may be exploited in the identification of their viewpoints, among other aspects associated to this research topic. After that, a case study is conducted, involving documents about brazilian politics written in 2010, year of brazilian presidential election. They are classified according to the \emph{pro-situation} and \emph{pro-opposition} viewpoints, and the words that are more emphasized by each viewpoint are also discussed.     

%A classificação de documentos de acordo com seus pontos de vista sobre um tema é um tópico de pesquisa que vem recebendo atenção crescente. Dado um conjunto de documentos sobre a política dos Estados Unidos, por exemplo, esse tipo de classificação pode ser empregado 
%para identificar quais deles defendem o ponto de vista \emph{liberal} e quais defendem o ponto de vista 
%\emph{conservador}. Esse tipo de classificação colabora com a análise de documentos argumentativos, especialmente em cenários nos quais não é possível lê-los um a um. Devido ao caráter desses documentos, esse tipo de classificação pode fazer parte de estudos interdisciplinares, envolvendo áreas como Ciência Política e Comunicação. 

%Diante disso, essa monografia apresenta uma revisão de trabalhos que tratam desse tópico, indicando quais são os principais classificadores utilizados, quais aspectos dos documentos podem ser explorados na identificação de seus pontos de vista, dentre outros aspectos associados a esse tópico de pesquisa. Em seguida, é feito um estudo de caso envolvendo documentos sobre a política brasileira escritos em 2010, ano de eleição presidencial. Eles são classificados de acordo com os pontos de vista \emph{pró-oposição} e \emph{pró-situação}, e as palavras mais enfatizadas por cada ponto de vista também são discutidas. 

\textbf{Keywords:} 
Opinion mining,
classification,
brazilian politics,
viewpoint.
\end{abstract}

%% As listas a seguir sao opcionais:
\listadefiguras
\listadetabelas
%\listadesiglas
%\listadesimbolos

\sumario

% O conteudo da monografia esta' nos seguintes arquivos:
\chapter{Introdução}

%\section{Motivação}

A busca por opiniões sempre desempenhou um papel importante na geração de novas escolhas. Antes de optar por assistir a um filme, é comum ler críticas a seu respeito ou considerar os comentários de outras pessoas; antes de comprar um produto, muitas vezes procuramos relatos sobre a satisfação de outros consumidores. Com a disseminação da Web e da Internet, a geração de opiniões com impacto, sobre os mais diversos assuntos, foi finalmente democratizada: não é mais preciso, por exemplo, ser um especialista em Economia ou Ciência Política para manter um \emph{blog} convincente sobre algum candidato às eleições.

Neste contexto, a busca por opiniões e comentários em \emph{sites}, \emph{blogs}, fóruns e redes sociais também se popularizou, passando a fazer parte do cotidiano dos consumidores \emph{online}. Uma pesquisa feita nos Estados Unidos revelou que entre 73\% e 87\% dos leitores de resenhas de serviços online, como críticas de restaurantes e albergues, sentem-se fortemente influenciados a consumi-los ou não a depender das opiniões contidas nessas resenhas \cite{pesquisa-eua}. Diante da relevância que opiniões têm na geração de decisões e no processo de consumo, estudos com o intuito de extraí-las da Web e interpretá-las automaticamente tornaram-se mais frequentes na área de Ciência da Computação. Juntos, esses estudos compõem o que ficou conhecido como \textbf{Análise de Sentimento} ou \textbf{Mineração de Opinião}, termos utilizados como sinônimos \cite{omsa} \cite{bingliu}.

De acordo com a \emph{survey} de Pang e Lee, referência mais completa para o estudo da área encontrada até o momento, a Mineração de Opinião envolve o emprego de diversas técnicas computacionais com o intuito de explorar algum dos tópicos abaixo \cite{omsa}:

\begin{enumerate}
    \item \textbf{Polaridade de sentimento ou graus de polaridade -} Dado um documento opinativo, para o qual se assume que as opiniões se referem basicamente a um único assunto, classifique-o como positivo ou negativo em relação a esse assunto (polaridades opostas) ou localize-o no espectro estabelecido entre duas polaridades opostas;
    \item \textbf{Detecção de subjetividade e identificação de opinião -} Dado um documento, detecte se ele é subjetivo ou não, constituindo-se de fatos ou opiniões, ou que partes dele são subjetivas;
    \item \textbf{Análise de tópico-sentimento -} Dado um documento opinativo, assume-se que suas opiniões podem se referir a tópicos diferentes, e deve-se identificar quais opiniões interagem com quais tópicos;
   \item \textbf{Pontos de vista ou perspectivas -} Dado um documento opinativo, que apresente uma perspectiva sobre um tema (um ponto de vista, uma orientação ideológica, um posicionamento), em vez de um sentimento polarizado sobre um único assunto, classifique-o de acordo com essa perspectiva;
   \item \textbf{Outras informações não-factuais -} Dado um documento com caráter emotivo/sentimental, identifique que tipos de humor o permeiam e/ou classifique-o de acordo com as emoções encontradas.

\end{enumerate}


%No começo dos anos 2000, os termos \emph{Opinion Mining} e \emph{Sentiment Analysis} são empregados pela primeira vez, referindo-se ao emprego de técnicas computacionais para identificação e avaliação de opiniões e sentimentos em textos \cite{panglee}. Esses termos têm funcionado como sinônimos \cite{bing}, aparecendo paralelamente na literatura. De acordo com \cite{panglee}, o termo \emph{Sentiment Analysis} é empregado pela primeira vez em 2001, nos artigos \textbf{Como citar referências citadas em outro trabalho?}; o termo \emph{Opinion Mining}, em 2003, no artigo \textbf{mesma duvida}. Nesta monografia, os dois termos serão utilizados de forma intercambiável.


A ideia de utilizar metodologias computacionais para identificar e analisar opiniões é muito anterior à popularização da Web \textbf{Citar artigos do fim da década de 60 e começo de 70 que provam isso}. Motivos: pouco dado, IR e ML imaturas. Explicar os 3 e como se relacionam com Natural Language Processing.
%\textbf{POPULARIDADE DAS ÁREAS DE ML E IR, SURGIMENTO DE DATASETS DISPONIVEIS POR CONTA DA POPULARIDADE DA INTERNET...}.

%One of the main reasons for the lack of study on opinions is that there was little opinionated text before
%the World Wide Web. Before the Web, when an individual needs to make a decision, he/she typically
%asks for opinions from friends and families. 

%Falar das principais tasks da área.

\section{Proposta}

Falar de Mineração de Perspectiva. Definir todos os termos correlatos utilizados, fechar os problemas da área e explicar como isso se diferencia de Opinion Mining clássica, que é basicamente Análise de Polaridade.

%\textbf{explicar que aqui muitas vezes é mais interessante entender posicionamentos globais do que a opinião sobre um único assunto, tópico, etc, e que tb há pesquisas na área de opinion mining que tratam disso; mostrar onde diferem, os problemas} Especificamente no que diz respeito à política, uma pesquisa realizada nos Estados Unidos \cite{rainiehorrigan}, com mais de 2500 usuários de Internet que compartilharam informação sobre as eleições de 2006 \emph{online}, concluiu o seguinte:

%\begin{enumerate}
 % \item Para 28\% deles, um dos principais motivos para realizar esta troca de informações \emph{online} era compreender as perspectivas de suas próprias comunidades; 
  %\item 34\% deles afirmaram que um dos principais motivos para realizar esta troca de informações \emph{online} era compreender as perspectivas exteriores à suas comunidades; 
   %\item 27\% deles pesquisaram avaliações políticas de organizações externas \emph{online};
%  \item 28\% deles afirmaram que a maioria dos \emph{sites} de política que frequentam compartilham com seus pontos de vista;
 % \item 29\% deles afirmaram que a maioria dos \emph{sites} de política que frequentam apresentam pontos de vista que divergem dos seus - indicação de que muitas pessoas não estão buscando apenas uma validação de opiniões pré-existentes\cite{panglee};
  %\item 8\% deles publicam seus próprios comentários políticos \emph{online}.
%\end{enumerate}  

%\textbf{Concluir fechando mais o escopo para perspectivas; tentar motivar um pouco para eleições, MAAAAAS falar brevemente de trabalhos que não necessariamente tem a ver com política e que tb tentam capturar perspectivas/pontos de vista}.

%Falar que as bases abrangentes serão discutidas, mas essa mono fecha o escopo em viewpoints, blablabla. Falar que se pretende fazer uma revisão comparativa dos trabalhos, mostrar alguma coisa prática na área com as decisões de projeto devidamente justificadas.

\section{Estrutura da Monografia}

Falar da metodologia de busca dos artigos 
%Modelo de monografia usando as normas ABNT (Associação Brasileira de Normas
%%Técnicas) \sigla{ABNT}{Associação Brasileira de Normas Técnicas}
%e adaptação personalizada 
%do padrão do Departamento de Ciência da Computação (DCC) da Universidade
%Federal da Bahia (UFBA).
%\sigla{DCC}{Departamento de Ciência da Computação}
%\sigla{UFBA}{Universidade Federal da Bahia}
%Fontes latex cedidos pela ABNT e disponibilizados por 
%Maurício Vieira. (Valeu Maurix!). Adaptado por Abelmon Bastos por solicitação
%da \profa\ Débora Abdalla para o semestre 2005.1.
%Adaptado por Rodrigo Rocha por solicitação da \profa\ Débora Abdalla no fim
%do semestre 2007.1.

% use sua propria estrutura

%Problema etc

%Objetivo etc

%Resultados esperados etc

%Estrutura/Organização etc

\chapter{Técnicas básicas e ferramentas utilizadas}

\textbf{Introduzir o capítulo quando tudo já estiver escrito.}
%Neste capítulo, serão descritas técnicas básicas de Aprendizado de Máquina e Processamento de Linguagem Natural, importantes para a compreensão dos capítulos seguintes desta monografia. Na seção \textbf{Classificadores}, serão discutidos os funcionamentos dos classificadores Naïve Bayes e \emph{Support Vector Machines} (SVMs), comuns em métodos de Mineração de Perspectiva baseados em frequências de palavras. Na seção \textbf{Modelos Gráficos}, serão discutidas a notação e hipóteses deste tipo de modelo, os \textbf{Modelos Generativos}, uma de suas subcategorias bastante explorada nesta monografia, e dois modelos generativos que, a partir da interpretação de documentos como uma mistura de tópicos, geram agrupamentos de palavras que facilitam a compreensão segmentada do conteúdo de um corpus. Estes modelos são o \emph{Latent Dirichlet Allocation} (LDA) e o \emph{Labeled-Latent Dirichlet Allocation} (L-LDA), utilizados em experimentos conduzidos ao longo de todo o projeto.

%\textbf{serão discutidos os modelos de tópicos \emph{Latent Dirichlet Allocation} (\textbf{LDA}) e \emph{Labeled Latent Dirichlet Allocation} (\textbf{L-LDA}). O segundo, que consiste em uma pequena modificação do primeiro, é utilizado em parte dos experimentos conduzidos neste projeto.}


\section{Modelos Gráficos}

%\begin{figure}[t]
%  \centering % este comando é usado para centralizar a Figura
%  \includegraphics[width=5cm, height=4cm]{graphical-naive.png}\\
%  \caption{Modelo gráfico do Naïve Bayes.}
%  \label{fig:graphical-naive}
%\end{figure}


\begin{figure}[t]
  \centering % este comando é usado para centralizar a Figura
  \includegraphics[width=5cm, height=3.4cm]{exemplo-modelo-grafico.png}\\
  \caption{Exemplo de modelo gráfico.}
  \label{exemplo:grafico}
\end{figure}



%\begin{figure}[t]
%  \centering % este comando é usado para centralizar a Figura
%  \includegraphics[width=5cm, height=4cm]{generative-bishop.png}\\
%  \caption{Modelo gráfico que representa o processo pelo qual objetos são gerados probabilisticamente. A distribuição de probabilidade para a imagem depende da identidade do objeto, de sua orientação e de sua posição \cite{bishop}.}
%  \label{fig:generative-bishop}
%\end{figure}

Modelos gráficos consistem na representação, através de um grafo, das relações entre um conjunto finito de variáveis aleatórias, provendo uma maneira simples de se representar distribuições de probabilidade \cite{bishop}. Cada vértice do grafo corresponde a uma variável aleatória (ou a um conjunto de variáveis aleatórias) ou a um parâmetro do modelo, e cada aresta reflete a relação entre dois vértices. Modelos gráficos são categorizados como dirigidos ou não-dirigidos. Por questões de escopo, apenas modelos dirigidos (também conhecidos como Redes Bayesianas) serão discutidos nesta seção.

Em um modelo gráfico dirigido, tem-se um grafo direcionado acíclico que representa a distribuição de probabilidade conjunta\footnote{Nesta monografia, os termos "distribuição de probabilidade conjunta" e "distribuição conjunta" serão utilizados de forma intercambiável.} para suas variáveis aleatórias. Cada aresta corresponde a uma distribuição de probabilidade condicional, incidindo no vértice cuja distribuição de probabilidade está condicionada ao valor do vértice de onde ela parte. Quando mais de uma variável tem distribuição de probabilidade condicionada aos mesmos vértices, é possível sintetizar a notação representando todas elas com um único vértice. Neste caso, o vértice fica dentro de um retângulo rotulado com o número de variáveis que ele representa. 

Por fim, vértices representados com círculos brancos correspondem a variáveis latentes - ou seja, cujos valores não são observáveis diretamente no conjunto de dados ao qual o modelo é aplicado; círculos cinzas, por sua vez, correspondem a variáveis observáveis, cujos valores estão explícitos no conjunto de dados. Variáveis latentes permitem que distribuições de probabilidade muito complexas, envolvendo variáveis observáveis, sejam construídas a partir de distribuições condicionais mais simples \cite{bishop}. A Figura \ref{exemplo:grafico} corresponde a um modelo gráfico com a seguinte distribuição conjunta

\begin{equation}
\label{eq:fiction}
\ensuremath{P(x | \alpha)P(y | x)\prod_{i=1}^{n}P(z_i | y, \beta)} 
\end{equation}

\ensuremath{\alpha} e \ensuremath{\beta} são parâmetros de distribuições de probabilidade, \ensuremath{x} e \ensuremath{y} são variáveis latentes e \ensuremath{z_1, ..., z_n}, sintetizadas na Figura \ref{exemplo:grafico} através do vértice \ensuremath{z}, são variáveis observáveis.


%A \textbf{Joint Distribution} do classificador Naïve Bayes, que corresponde ao produto de todas as distribuições de probabilidade do modelo, pode ser escrita como 

% rótulo indica o número de variáveis. % esta represehaja, no modelo, um certo número de variáveis aleatórias que dependem 

 

%\textbf{********** Pensar num jeito de explicar a notação sem falar de Naïve Bayes. Falar de variáveis latentes e observáveis. *****************}

%\textbf{                                                                                The
%primary role of the latent variables is to allow a complicated distribution over the
%observed variables to be represented in terms of a model constructed from simpler
%(typically exponential family) conditional distributions.}

%\subsubsection{Modelos Generativos} 

Uma categoria de modelos gráficos explorada neste projeto são os \textbf{modelos generativos}. Eles associam distribuições de probabilidade a todas as variáveis aleatórias envolvidas, permitindo a geração - i.e. simulação - de seus valores. Modelos generativos são úteis para expressar os processos pelos quais dados observáveis são obtidos. Em um modelo deste tipo, os valores das variáveis aleatórias podem ser obtidos através de técnicas de amostragem aplicadas à distribuição de probabilidade conjunta. Além do Naïve Bayes, discutido na seção \ref{subsection:naive}, os modelos gráficos generativos \emph{Latent Dirichlet Allocation} (LDA) e \emph{Labeled Latent Dirichlet Allocation} (L-LDA) foram empregados em experimentos ao longo de todo o projeto. Uma discussão sobre eles pode ser encontrada na seção \ref{subsection:LDA}. %, como no seguinte exemplo: dado um conjunto de imagens de objetos, variáveis latentes podem ser utilizadas para representar suas posições e orientações; em seguida, para identificar o objeto contido em uma imagem particular observada, deve-se obter a distribuição de probabilidade a posteriori para objetos, o que envolve integrar sobre todas as posições e orientações possíveis \cite{bishop}. A Figura \textbf{2.2} indica as distribuições de probabilidade envolvendo variáveis aleatórias para posição, orientação, imagem e objeto.  %, como será discutido, no caso particular do modelo Naïve Bayes, na seção \ref{subsection:naive}.

%, como a amostragem de Gibbs apresentada na seção 2.1.1. É um mecanismo diferente do empregado em SVMs, por exemplo, em que variáveis observáveis não são associadas a nenhuma distribuição de probabilidade e as categorias são estimadas diretamente, condicionadas a seus valores. SVMs estão apresentados com mais detalhes na seção 2.1.2. \textbf{Acho q a monografia n precisa falar de modelos discriminativos.}

 %  modelo de tópicos L-LDA, que consiste em uma alteração simples no modelo LDA. Estes últimos, discutidos na seção \ref{subsection:LDA}, %, bem como o classificador Naïve Bayes.

\subsection{Naïve Bayes}
\label{subsection:naive}

O Naïve Bayes é um modelo gráfico generativo que assume a independência condicional das características \ensuremath{F_1, ..., F_k} presentes em um conjunto de documentos \ensuremath{D}. Se \ensuremath{D} é composto de documentos de texto, estas características normalmente correspondem a todas as suas palavras distintas. Isto equivale a assumir, portanto, que a presença de uma palavra em um documento qualquer não é informativa sobre a presença de nenhuma outra. 

A finalidade básica do Naïve Bayes é estimar a probabilidade de um documento \ensuremath{d} pertencer a uma certa classe \ensuremath{c}. Para isto, \ensuremath{d} é representado de forma simplificada, através de um vetor \ensuremath{v_d} em que cada posição corresponde a uma de suas \ensuremath{n} palavras. Com esta representação, a probabilidade de \ensuremath{d} pertencer a uma classe \ensuremath{c} pode ser calculada via Teorema de Bayes como

%variáveis aleatórias, em que \ensuremath{d_i denota a relevância de F_i para d}. \textbf{MEXER NESSA EXPLICAÇÃO!} \ensuremath{d_i} pode ser a frequência de \ensuremath{F_i} em \ensuremath{d}, ou indicar sua presença com o valor 1 e sua ausência com o valor 0, por exemplo. 

\begin{equation}
\label{eq1:naive}
\ensuremath{P(c | v_{d1}, ..., v_{dn}) = \frac{p(c) \times p(v_{d1}, ..., v_{dn} | c)}{p(F_1, ..., F_k)} }
\end{equation}

Como o Naïve Bayes assume que as palavras dos documentos são condicionalmente independentes, a equação \ref{eq1:naive} pode ser reescrita como 

\begin{equation}
\label{eq2:naive}
\ensuremath{P(c | v_{d1}, ..., v_{dn}) = \frac{p(c) \times p(v_{d1} |  c)p(v_{d2} | c)...p(v_{dn - 1} | c)p(v_{dn} | c)}{p(F_1, ..., F_k)} = \frac{p(c) \times \prod_{i = 1}^{n} p(v_{di} |  c)}{p(F_1, ..., F_k)}}
\end{equation}

%Na \textbf{seção XXX}, a equação \ref{eq2:naive} será retomada para a discussão sobre como construir um classificador a partir do modelo Naïve Bayes.

Como o Naïve Bayes é um modelo generativo, ele permite que se simule a criação de um documento \ensuremath{d}, pertencente a uma classe \ensuremath{c}, através da amostragem de suas variáveis aleatórias. Sem perda de generalidade, assume-se que \ensuremath{c} é uma variável aleatória que pode assumir \ensuremath{m} valores naturais distintos, variando de 0 a \ensuremath{m - 1}. Cada valor corresponde a uma classe diferente, sendo escolhido de acordo com  

\begin{equation}
\label{eq3:label-naive}
\ensuremath{c \sim Binomial(m - 1, \pi)}
\end{equation}

Antes de iniciar o processo de geração de documentos, define-se um parâmetro \ensuremath{\pi} para a distribuição binomial em \ref{eq3:label-naive} de acordo com

\begin{equation}
\label{eq4:pi-naive}
\ensuremath{\pi \sim Beta(\alpha, \beta)}
\end{equation}

\ensuremath{\alpha} e \ensuremath{\beta} são denominados \emph{hiperparâmetros}, pois são parâmetros de uma distribuição através da qual se escolhe um dos parâmetros do modelo - no caso, \ensuremath{\pi} \cite{gibbs}. Após uma classe ter sido fixada de acordo com \ref{eq3:label-naive}, seleciona-se uma palavra para cada posição \ensuremath{j} do vetor \ensuremath{v_d}, de acordo com uma distribuição de probabilidade sobre \ensuremath{F_1, ..., F_k}

\begin{equation}
\label{eq5:multi-naive}
\ensuremath{v_{dj} \sim Multinomial(F_1, ..., F_k, \theta_c)}
\end{equation}

A distribuição utilizada depende do valor de \ensuremath{c} amostrado anteriormente, de modo que há \ensuremath{m} parâmetros \ensuremath{\theta_c}. Cada \ensuremath{\theta_c} é escolhido antes do processo de geração dos documentos, de acordo com

\begin{equation}
\label{eq6:theta-naive}
\ensuremath{\theta_c \sim Dirichlet(\gamma_c)}
\end{equation}

\ensuremath{\gamma_c}, assim como \ensuremath{\alpha} e \ensuremath{\beta}, é um hiperparâmetro.

A distribuição conjunta para este modelo generativo é dada por

\begin{equation}
\label{eq7:joint-naive}
\ensuremath{P(\pi | \alpha, \beta)\prod_{i = 0}^{m - 1}P(\theta_{c = i} | \gamma_{c=i})\prod_{j = 1}^{|D|}P(c_j | \pi)P(v_d^{(j)} | \theta_{c_j}, c_j)}
%\prod_{i = 0}^{m - 1}P(c = i | \pi)P(\theta_{c = i} | \gamma_{c = i})P(v_d^{(j)} | \theta_{ c = i})} 
\end{equation}

\begin{figure}[t]
  \centering % este comando é usado para centralizar a Figura
  \includegraphics[width=5cm, height=4.5cm]{naive-joint.png}\\
  \caption{Modelo gráfico Naïve Bayes.}
  \label{naive:joint}
\end{figure}


em que \ensuremath{c_j} é a classe selecionada para o \ensuremath{j}-ésimo documento de \ensuremath{D} e \ensuremath{v_d^{(j)}} é seu vetor de palavras. A Figura \ref{naive:joint} representa esta distribuição conjunta de forma gráfica, com vértices para variáveis aleatórias e relações de probabilidade condicional evidenciadas pelas arestas.


% Caso haja \ensuremath{m} classes, tem-se que a classe de \ensuremath{d
 % Assume-se que há uma variável aleatória para cada \ensuremath{f \in F}, cujos valores denotam a relevância de cada uma delas para um \ensuremath{d} fixado. %Estas variáveis podem assumir valores discretos, como a contagem de quantas vezes elas ocorreram em \ensuremath{d}os valores 1 e 0, caso façam parte ou não do documento \ensuremath{d} e 0 em caso contrário, ou contí relevância pode ser medida binariamente - \ensuremath{1 se f é uma das palavras de d, 0 em caso contrário} -  %, e a presença ou ausência de uma delas em um documento qualquer é independente da presença ou ausência de qualquer outra. 

% parte da ideia de que as informações presentes em um documentoe, são independentes entre si. No caso de um documento de texto, assume-se que 




%associa uma classe a cada documento de um corpus. 

% No caso do classificador Naïve Bayes - modelo gráfico que, por conveniência, é discutido mais detalhadamente na seção \ref{subsection:bayes} -, tem-se variáveis aleatórias que representam a classe de cada documento e, para todos os documentos, variáveis aleatórias para cada um de seus termos distintos. Este modelo contém dois parâmetros, \ensuremath{\pi} e \ensuremath{\theta}, que ajustam as distribuições de probabilidade das variáveis aleatórias. 


%Este tipo de classificador parte da ideia de que as informações presentes em um documento, utilizadas na determinação de sua classe, são independentes entre si. No caso de um documento de texto, assume-se que a presença ou ausência de um termo - uma palavra ou uma sequência de palavras - é independente da presença ou ausência de qualquer outro. Definida esta hipótese, a probabilidade de que um documento \emph{d} seja de uma classe \emph{c} é tal que

%\begin{equation}
%label{eq1}
%\ensuremath{P(c|d) \propto P(c)\prod_{k=1}^{t_d}P(t_k|c)}  
%\end{equation} 

%onde \ensuremath{P(t_k|c)} é a probabilidade condicional do termo \ensuremath{t_k}  ocorrer em um documento da classe \ensuremath{c},  \ensuremath{P(c)} é a probabilidade a priori de um documento qualquer pertencer à classe \ensuremath{c} e \ensuremath{t_d} é o número de termos em \emph{d} \cite{stanford-IRbook}.

%As probabilidades envolvidas na relação \ref{eq1} contêm integrais díficeis ou mesmo impossíveis de se calcular analiticamente. Para calcular \ensuremath{P(c|d)}, portanto, utiliza-se aproximações obtidas através de técnicas de amostragem. Uma destas técnicas, comum na literatura de Aprendizado de Máquina e empregada neste projeto, é a amostragem de Gibbs. Em uma iteração da amostragem, a técnica condiciona as probabilidades calculadas para um documento \ensuremath{k} às classificações obtidas para os \ensuremath{k - 1} documentos anteriores \cite{resnik-gibbs}. A técnica está descrita, de forma básica, no algoritmo \textbf{X (o algorithm2e.sty tá dando pau)}.
                             % The basic idea in Gibbs sampling is that, rather than probabilistically picking
%the next state all at once, you make a separate probabilistic choice for each of the k dimensions, where each
%choice depends on the other k − 1 dimensions.


%A hipótese de independência entre os termos, razão pela qual o Naïve Bayes tem este nome\footnote{Naïve é uma palavra de origem francesa que significa "ingênua"}, simplifica bastante a estrutura da informação contida nos documentos. Ainda assim, o classificador costuma apresentar boas performances em categorização de textos, sendo utilizado, por exemplo, como base metodológica para alguns filtros de \emph{spam} \cite{paul}. Para melhorar o desempenho do Naïve Bayes, é comum fixar um conjunto de documentos previamente classificados de forma correta e utilizar a informação sobre suas classes na determinação das classes de outros documentos. Ao conjunto de documentos previamente classificados, dá-se o nome de \textbf{conjunto de treinamento}; ao conjunto de documentos a serem classificados, \textbf{conjunto de teste}.

%Todos os experimentos com Naïve Bayes conduzidos neste projeto utilizam a implementação disponível no repositório \emph{online} de Aline Bessa \cite{alibezz-nb}. O número de iterações para a Amostragem de Gibbs foi fixado em 500. %Através do uso de Amostragem de Gibbs, a cada iteração se obtém uma aproximação melhor O número de amostras coletadas envolvendo cada classe foi fixado em 500. 



%A Figura 2.1, correspondente ao classificador Naïve Bayes, contém toda a notação necessária para as discussões sobre modelos gráficos deste projeto. Assume-se que o classificador será aplicado a um conjunto de documentos \ensuremath{D}. A classe de cada documento \ensuremath{j}, \ensuremath{C_j}, é uma variável aleatória com distribuição de probabilidade \ensuremath{P(C_j|\pi)}. A relação entre \ensuremath{C_j} e \ensuremath{\pi}, evidenciada em \ensuremath{P(C_j|\pi)}, é representada na Figura 2.1 pela aresta que sai de \ensuremath{\pi} e incide em \ensuremath{C_j}. \ensuremath{C_j} está no interior de um retângulo rotulado com \ensuremath{|D|}, o que significa que há \ensuremath{|D|} variáveis deste tipo, uma para cada documento. %, com \ensuremath{j} variando de 1 a \ensuremath{N}. 

%Cada documento \ensuremath{j}, por sua vez, contém um conjunto de \ensuremath{T_j} termos distintos. A probabilidade associada a cada termo \ensuremath{T_{jk}}, \ensuremath{k} variando de 1 a \ensuremath{T_j}, é condicionada ao valor da variável \ensuremath{C_j}. Desse modo, e considerando o parâmetro \ensuremath{\theta}, tem-se, por documento, \ensuremath{T_j} distribuições de probabilidade do tipo \ensuremath{P(T_{jk}|C_j, \theta)}. Uma vez mais, o rótulo do retângulo mais interno indica o número de variáveis, sintetizadas por um único círculo, que estabelecem o mesmo tipo de relação com \ensuremath{\theta} e \ensuremath{C_j}. %no modelo. 

%Por fim, as variáveis representadas por círculos brancos são latentes - ou seja, seus valores não são observáveis diretamente no conjunto de dados ao qual o modelo é aplicado. As variáveis em cinza, por sua vez, correspondem a dados diretamente observáveis; neste caso, às palavras distintas de cada documento. A \textbf{Joint Distribution} do classificador Naïve Bayes, que corresponde ao produto de todas as distribuições de probabilidade do modelo, pode ser escrita como 

%\begin{equation}
%\label{generative:eq1}
%\ensuremath{\prod_{j=1}^{|D|} \bigg\{P(C_j|\pi)\bigg[\prod_{k=1}^{T_j}P(T_{jk}|C_j)\bigg]\bigg\}}  
%\end{equation} 

%Adequando-se a notação, observa-se que a relação \ref{eq1}, apresentada na seção 2.1.1, pode ser facilmente obtida via a \textbf{Joint Distribution} \ref{generative:eq1}.

%\subsubsection{Modelos Generativos} 

%Uma categoria de modelos gráficos explorada neste projeto são os Modelos Generativos. Modelos deste tipo associam distribuições de probabilidade a todas as variáveis aleatórias envolvidas, permitindo a geração - i.e. simulação - de seus valores. Modelos generativos são úteis para expressar os processos pelos quais dados observáveis são obtidos, como no seguinte exemplo: dado um conjunto de imagens de objetos, variáveis latentes podem ser utilizadas para representar suas posições e orientações; em seguida, para identificar o objeto contido em uma imagem particular observada, é preciso obter a distribuição de probabilidade a posteriori para objetos, o que envolve integrar sobre todas as posições e orientações possíveis \cite{bishop}. A Figura 2.2 indica as distribuições de probabilidade envolvendo variáveis aleatórias para posição, orientação, imagem e objeto.

%Com um modelo generativo, os valores das variáveis podem ser obtidos através de técnicas de amostragem aplicadas à \textbf{Joint Distribution}, como a amostragem de Gibbs apresentada na seção 2.1.1. É um mecanismo diferente do empregado em SVMs, por exemplo, em que variáveis observáveis não são associadas a nenhuma distribuição de probabilidade e as categorias são estimadas diretamente, condicionadas a seus valores. SVMs estão apresentados com mais detalhes na seção 2.1.2. \textbf{Acho q a monografia n precisa falar de modelos discriminativos.}

%Dois modelos generativos são bastante explorados nos capítulos seguintes desta monografia: o classificador Naïve Bayes e o modelo de tópicos L-LDA, que consiste em uma alteração simples no modelo LDA. Estes últimos, discutidos na seção \ref{subsection:LDA}, foram empregados em experimentos e estudos de caso ao longo de todo o projeto. %, bem como o classificador Naïve Bayes.


% são mapeados em categorias baseando-se nesta informação condicionando 


%os valores das variáveis latentes são associados a distribuições de probabilidade que as condicionam a variáveis observ condicionados a variáveis observadas explicitamente.

%por exemplo, em que o valor de uma variável aleatória


%Discriminative models differ from generative models in that they do not allow one to generate samples from the joint distribution of x and y.
% Quando uma técnica de amostragem adequada é aplicada ao modelo, a fim de obter valores para suas variáveis, tem-se a simulação de dados cuja distribuição de probabilidade aproxima-se daquela dos dados 

%\textbf{Figura do bayes}
%\textbf{Figura bishop com legenda}    

%=> Falar de modelos generativos e fazer link com o resto da seção.
% A independência condicional \pi é um hiperparâmetro do modelo, e a caixa mais externa, rotulada com um \ensuremath{N}, indica que \ensuremath{j} varia de 1 a \a relação entre \ensuremath{\pi} e as variáveis   

\subsection{LDA}
\label{subsection:LDA}


O modelo LDA parte da ideia de que um documento pode tratar de múltiplos tópicos, refletidos nas palavras que o compõem \cite{pnas}. Assim como no modelo Naïve Bayes, palavras podem ser geradas de acordo com distribuições multinomiais específicas. A diferença é que, no LDA, cada palavra é gerada a partir de uma mistura de tópicos; no Naïve Bayes, elas são geradas a partir de um só tópico (classe) \cite{gibbs-lingpipe}.

O LDA associa as palavras dos documentos a tópicos diferentes, com maior ou menor probabilidade, criando agrupamentos que se relacionam semanticamente. Os tópicos em um LDA são variáveis latentes, cujos significados requerem uma interpretação posterior ao processamento. Esta interpretação baseia-se nas relações semânticas entre as palavras que se associaram mais fortemente a cada um deles. %, algo relativamente subjetivo. %ser bastante levando  

Para ilustrar como as palavras evidenciam o significado de um tópico, um experimento envolvendo receitas culinárias extraídas do \emph{site} \textbf{allrecipes.com} foi executado. Apenas os ingredientes de cada receita foram considerados. Na Tabela 2.1, constam as cinco palavras mais fortemente associadas a quatro tópicos, de acordo com o LDA. % - ou seja, as cinco palavras para as quais as probabilidades obtidas com a equação \ref{eq2} foram mais altas. %que obtiveram mais alta probabilidade de acordo com \ref{eq2}.

\begin{table}[t]
\centering
\label{LDA:1}
\begin{tabular}{| l | p{7cm} | }
\hline
\textbf{Tópico} & \textbf{Palavras} \\ \hline
\textbf{Tópico 1} & beef, cheese, tomato, sauce, pepper \\ \hline
\textbf{Tópico 2} & chicken, breast, pastum, broth, tomato \\ \hline
\textbf{Tópico 3} & flour, sugar, butter, powder, egg  \\ \hline
\textbf{Tópico 4} & cream, cheese, butter, milk, cake \\ \hline
\end{tabular}
\caption{As cinco palavras mais fortemente associadas a quatro tópicos gerados por um LDA.}
\end{table}

Considerando que todas as receitas pertencem à culinária tradicional dos Estados Unidos, as palavras listadas na Tabela 2.1, e a forma como se associam em torno de cada tópico, são indicativos do bom funcionamento do LDA. O primeiro tópico pode ser interpretado como \textbf{ingredientes para          \emph{cheeseburger}}; o segundo, ao associar \emph{chicken}, \emph{pastum} e \emph{broth}, remete a receitas de sopas e caldos comuns em climas frios, podendo ser interpretado como \textbf{ingredientes para sopa}; o terceiro pode ser interpretado como \textbf{ingredientes para bolo}; o quarto, ao associar \emph{cream}, \emph{cheese} e \emph{cake}, pode ser interpretado como \textbf{ingredientes para \emph{cheesecake}}. É importante frisar que estas interpretações, apesar de subjetivas, indicam perspectivas culinárias distintas e coerentes internamente. Seria diferente de encontrar, por exemplo, um tópico fortemente associado às palavras \emph{sugar}, \emph{pepper} e \emph{potato}, dificilmente encontradas em uma mesma receita típica dos Estados Unidos.


%- ou seja, não são identificados antes ou durante o processamento -cujos significados requerem  só podem ser interpretados após o processamento do modelo, observando-se  

% que, após o processamento, sintetizam um significado associado ao


%atribuir um significado a cada um deles apenas observando as relações entre as palavras associadas com mais destaque. Para o LDA, os tópicos não são, portanto, pré-identificados antes do processamento, requerendo uma interpretação posterior e subjetiva de seus significados. %Neste sentido, o modelo é útil para identificar padrões contidos em documentos de texto.  

O modelo LDA trata cada documento pertencente a um conjunto de documentos \ensuremath{D} como uma mistura de tópicos, representada por uma distribuição de probabilidade sobre um conjunto de tópicos \ensuremath{T}. Cada tópico, por sua vez, é visto como uma mistura de palavras, representada por uma distribuição sobre todas as palavras distintas de \ensuremath{D}. Para cada documento \ensuremath{d \in D}, é fixada uma distribuição de probabilidade sobre tópicos \ensuremath{\theta_d}, condicionada a um hiperparâmetro \ensuremath{\alpha}

\begin{equation}
\label{lda:topic}
\ensuremath{\theta_d \sim Dirichlet(\alpha)}
\end{equation}

Para cada tópico \ensuremath{t \in T}, é fixada uma distribuição de probabilidade \ensuremath{\phi_t} sobre palavras, condicionada a um hiperparâmatro \ensuremath{\beta}

\begin{equation}
\label{lda:word}
\ensuremath{\phi_t \sim Dirichlet(\beta)}
\end{equation}

Em seguida, para cada uma das \ensuremath{n} palavras de \ensuremath{d}, um tópico \ensuremath{t} é escolhido, de acordo com \ensuremath{\theta_d}

\begin{equation}
\label{lda:topic-chosen}
\ensuremath{t \sim Discrete(\theta_d)}
\end{equation}

e uma palavra \ensuremath{w} é gerada de acordo com \ensuremath{\phi_t}

\begin{equation}
\label{lda:word-chosen}
\ensuremath{w \sim Discrete(\phi_t)}
\end{equation}

A distribuição conjunta do modelo LDA é dada por

\begin{equation}
\label{joint:lda}
\ensuremath{\prod_{i=1}^{|D|} \bigg\{P(\theta_i|\alpha)\bigg[\prod_{j=1}^{|W|}P(t_j|\theta_i)P(w_j|t_j,\beta)\bigg]\bigg\}}  
\end{equation}

\ensuremath{P(w_j|t_j,\beta)} reflete o quanto a palavra \ensuremath{w_j} se relaciona com o tópico \ensuremath{t_j}; \ensuremath{P(t_j|\theta_i)}, por sua vez, funciona como uma medida do quanto o tópico \ensuremath{t_j} é importante no contexto do documento \ensuremath{d_i} \cite{pnas}. Na Figura \ref{fig:lda}, tem-se o modelo gráfico do LDA correspondente à distribuição conjunta \ref{joint:lda}. 

\begin{figure}[t]
  \centering % este comando é usado para centralizar a Figura
  \includegraphics[width=6.5cm, height=3.2cm]{Latent_Dirichlet_allocation.png}\\
  \caption{Modelo gráfico LDA.} %\ref{joint:lda}.}
  \label{fig:lda}
\end{figure}

A implementação do LDA utilizada para estes experimentos está disponível no repositório \emph{online} de Alexandre Passos \cite{top-lda}, e o número de iterações para amostragem de tópicos e palavras foi fixado em 100. 


%O cálculo destas probabilidades também envolve integrais difíceis, ou mesmo impossíveis, de resolver analiticamente, o que implica no uso de técnicas de amostragem e aproximação. A amostragem de Gibbs, descrita superficialmente na seção \ref{subsection:bayes}, é uma alternativa para a geração de valores para variáveis aleatórias contidas no modelo \cite{gibbs-lingpipe}, sendo utilizada nos experimentos com LDA conduzidos neste projeto. 

%Quando um método de classificação é aplicado a um conjunto de documentos, a interpretação do resultado é imediata: cada documento estará associado a uma única classe. No caso do modelo de tópicos LDA, a interpretação do resultado é mais subjetiva. É preciso observar as palavras que se associam com maior probabilidade a cada tópico, buscando algum tipo de semelhança entre elas, para inferir seus significados. 


% o agrupamento de palavras como \emph{beef}, \emph{sauce} e \emph{cheese} em torno de um mesmo tópico indica o bom funcionamento do LDA%As palavras associadas a cada tópico siderando que as receitas pertencem à culinária tradicional dos Estados Unidos. O primeiro tópico pode ser interpretado como \textbf{Receitas com Carne}, pois associa ingredientes comumente consumidos  
%\textbf{Análise}

\subsubsection{L-LDA}

O L-LDA é uma variação do LDA em que se restringe o número de tópicos associados a cada documento. Ou seja, as distribuições fixadas para os tópicos de cada documento não necessariamente são sobre todos os tópicos \ensuremath{t \in T}. Além disso, os tópicos presentes em cada documento são identificados antes da execução do modelo, o que diminui a subjetividade envolvida na interpretação de seus significados após o processamento. 

Um bom exemplo para ilustrar a aplicação deste modelo envolve um \emph{blog}, em que cada \emph{post} é marcado com um conjunto específico de \emph{tags}. Se cada \emph{tag} é interpretada como um tópico, é possível informar ao L-LDA em que \emph{posts} cada uma delas está presente, processar os \emph{posts} com o modelo e saber, após o processamento, quais palavras se associam mais fortemente a cada \emph{tag}. Neste exemplo, existe um mapeamento direto entre os tópicos e as \emph{tags}, conduzindo a uma interpretação mais imediata do significado de cada agrupamento de palavras.

Experimentos com o L-LDA foram desenvolvidos ao longo deste projeto, associando cada tópico, por exemplo, a uma perspectiva a ser minerada. Após a execução do modelo, as palavras mais fortemente associadas a cada perspectiva ilustram como os assuntos discutidos são enfocados por elas. Quanto mais duas perspectivas se distanciam, mais diferentes são as palavras que se associam com destaque a cada uma delas.

O L-LDA foi discutido pela primeira vez em um artigo de Ramage et al. \cite{llda}, aplicado ao problema de atribuição de crédito em páginas do \emph{site del.icio.us}, marcadas com múltiplas \emph{tags}. O artigo parte da hipótese de que, embora um documento possa estar marcado com várias \emph{tags} diferentes, nem sempre elas se aplicam igualmente a todas as palavras nele contidas. A ideia da atribuição de crédito, portanto, consiste em associar cada palavra do documento às \emph{tags} mais apropriadas e vice-versa. 

%Uma aplicação de L-LDA, sugerida neste artigo e bastante empregada neste projeto, envolve, após a associação de palavras a tópicos (neste caso, \emph{tags}), a extração de trechos do documento que as contenham. Isto conduz a uma compreensão melhor de como o conteúdo se associa aos tópicos separadamente. 

A implementação de L-LDA utilizada neste projeto também está disponível no repositório \emph{online} de Alexandre Passos \cite{top-llda}. O número de iterações para amostragem de tópicos e palavras, em todos os experimentos, foi fixado em 100.
% do envolvendo uma ou mais palavras associadas a uma \emph{tag} particular. 
%O resultado de um método de classificação consiste em associar cada documento de um conjunto Diferentemente de métodos de classificação, cujo resultado pode ser interpretado de forma obj

%Assim como discutido na seção \ref{subsection:bayes}, 
% tópicos é fixada multinomial  \ensure\ensuremath{d} se associa a tópicos \ensuremath{t \in T} com probabilidades diferentes, o que é determinado pela escolha de uma distribuição de probabilidade sobre os tópicos. Em seguida, cada tópico é tratado como uma distribuição sobre palavras e cada palavra \ensuremath{w_i} de \ensuremath{d} se associa a eles de acordo com estas distribuições. A probabilidade da \ensuremath{i}-ésima palavra de \ensuremath{d} pode ser calculada como 

%P(wi) = sum P(wi | ti = j) P(ti = j), j de 1 a T. 

%onde ti é uma variável latente - ou seja, cujo valor não é observável, mas sim inferido - referente a um tópico, P(wi | ti = j) é a probabilidade de wi se associar ao tópico j e P(ti = j) é a probabilidade de se obter o tópico j na distribuição fixada para o documento \ensuremath{d} \cite{pnas}.

 %\textbf{Falar das receitas.}
%\textbf{Falar do uso de Gibbs Sampling para inferência}


%\subsection{L-LDA}


\section{Classificadores}

A classificação de documentos de texto de acordo com suas perspectivas é um dos principais objetivos dos trabalhos revisados neste projeto. Grande parte deles utiliza os classificadores Naïve Bayes ou \emph{Support Vector Machines} (SVMs), apresentados respectivamente nas seções \ref{subsection:naive} e \ref{subsection:SVMs}, como parte de suas metodologias. O desempenho destes classificadores é comumente medido através das seguintes métricas: \textbf{taxa de acerto}, \textbf{precisão}, \textbf{rechamada} ou \textbf{métrica F1}.

A taxa de acerto é definida pela razão entre o número de documentos classificados corretamente e todos os documentos avaliados. A precisão, medida para uma classe \ensuremath{c} qualquer, é definida pela razão entre o número de documentos classificados corretamente como \ensuremath{c} e todos os documentos classificados como \ensuremath{c}. A rechamada, medida também para uma classe \ensuremath{c} qualquer, é definida pela razão entre o número de documentos classificados corretamente como \ensuremath{c} e a soma deste valor com o número de documentos classificados erroneamente para todas as demais classes. A métrica F1, também medida para uma classe \ensuremath{c} qualquer, é dada por

\begin{equation}
\ensuremath{2 \times \frac{precisao \times rechamada}{precisao + rechamada}}
\end{equation} 

Estas métricas revelam aspectos diferentes do desempenho de um classificador. Por este motivo, é comum encontrar mais de uma delas sendo utilizada no mesmo contexo. Além de medirem desempenho, elas estabelecem critérios objetivos para a comparação entre métodos de classificação, como pode ser visto nos artigos de Lin et al. sobre o conflito Israel-Palestina \cite{lin-et-al2006} e de Efron sobre orientação cultural \cite{efron}. 

\subsection{Naïve Bayes}
\label{subsection:bayes}

O modelo Naïve Bayes foi apresentado na seção \ref{subsection:naive}. Nesta seção, será discutido como construir um classificador de documentos a partir dele. Esta seção trata de documentos de texto em particular, por se tratarem do objeto básico de estudo deste projeto. 
Ra
Sabe-se que, em um Naïve Bayes, assume-se que as as características em um documento são condicionalmente independentes, o que equivale a afirmar, por exemplo, que a presença de uma palavra em um documento de texto não é informativa sobre a presença de nenhuma outra. Apesar desta hipótese simplificar bastante a estrutura linguística de um texto, classificadores construídos a partir do modelo Naïve Bayes, denominados classificadores Naïve Bayes, reportam um bom desempenho em várias tarefas de classificação baseadas em palavras \cite{naive-at-forty} \cite{mccallum-naive}.

Dados um documento \ensuremath{d} pertencente a um conjunto de documentos \ensuremath{D}, todas as palavras distintas de \ensuremath{D}, \ensuremath{F_1, ..., F_k}, uma variável aleatória \ensuremath{c}, que representa as possíveis classes de \ensuremath{d}, e um vetor \ensuremath{v_d}, em que cada posição corresponde a uma de suas \ensuremath{n} palavras, tem-se que
 
\begin{equation}
\label{eq2:bayes}
\ensuremath{P(c | v_{d1}, ..., v_{dn}) = \frac{p(c) \times \prod_{i = 1}^{n} p(v_{di} |  c)}{p(F_1, ..., F_k)}}
\end{equation}

conforme discutido anteriormente na seção \ref{subsection:naive}. Um classificador Naïve Bayes deve rotular o documento \ensuremath{d} com o valor de \ensuremath{c} que maximiza a equação \ref{eq2:bayes}. Como o denominador na equação \ref{eq2:bayes} é o mesmo para todas as classes, ele pode ser ignorado nestes cálculos.

Normalmente, classificadores Naïve Bayes são utilizados de forma semi-supervisionada. Isto significa que eles são submetidos a uma etapa de treinamento, na qual aprendem as classes associadas a alguns documentos, e a uma etapa de classificação, na qual devem simular o processo gerador destes documentos e utilizar esta informação para classificar outros. Basicamente, as informações aprendidas na etapa de treinamento modelam as distribuições das palavras de \ensuremath{D} por classe, gerando parâmetros para as distribuições de probabilidade envolvidas na classificação de outros documentos. 

Todos os experimentos com um classificador Naïve Bayes conduzidos neste projeto utilizam a implementação disponível no repositório \emph{online} de Aline Bessa \cite{alibezz-nb}. O número de iterações para a amostragem de documentos e classes, em todos os experimentos, foi fixado em 500.

%O número de iterações para amostragem de tópicos e palavras, em todos os experimentos, foi fixado em 100.

% a amostragem da equação \ref{eq:bayes} foi fixado em 500.

% classificando outros                        Dado que tenho exemplos de texto de cada classe. O
%que posso inferir sobre o processo gerador destes textos


%                              e                  o classificador reporta
%o melhor desempenho em v ́rias tarefas de classifica ̧ ̃o. Este fenˆmeno  ́
%                        a                         ca

 
%\begin{algorithm}
%\ensuremath{z^{(0)}} \gets \ensuremath{\langle z_1^{(0)}, \ldots, z_k^{(0)}\rangle}
%\FOR{\ensuremath{t = 1} to \ensuremath{T} do}
 % \FOR{\ensuremath{i = 1} to \ensuremath{k} do}
 %   \ensuremath{z_i^{(t + 1)} ~ P(Z_i | z_1^{(t + 1)}, \ldots, z_{i - 1}^{(t + 1)}, z_{i + 1}^{(t)}, \ldots, z_k^{(t)})}
 % \ENDFOR
%\ENDFOR

%\end{algorithm}


%Além disso, você deveria começar a descrever o naive bayes falando que
%ele é um modelo generativo e que assume independência condicional das
%features dado as classes. Dizer que naive bayes aproxima P(vetor(d),c)
%como sendo P(c) Produtorio(P(di|c), o que é a mesma coisa que dizer
%que as palavras são condicionalmente independentes dado as classes,
%mostrar como vc faz pra usar isso pra classificar, usando P(c|d)  =
%P(c,d)/P(d), e que como P(d) independe da classe pode ser ignorado se
%você quer encontrar a melhor classe. Daí você deveria especificar de
%onde vêm essas probabilidades, colocando uma prior Beta(alpha, alpha)
%em P(c) e uma prior Dirichlet(alpha) em P(tk|c).

%Dado um conjunto de documentos \ensuremath{D} e um conjunto de classes \ensuremath{C}, o classificador Naïve Bayes estima, através da aplicação do Teorema de Bayes,  a probabilidade de cada \ensuremath{d \in D} ser de cada uma das classes \ensuremath{c \in C}. Com estas probabilidades calculadas, o classificador determina, para todo \ensuremath{d}, qual é a classe \emph{c} a que ele estará associado. Esta classe pode, por exemplo, ser aquela para a qual a probabilidade obtida foi simplesmente a mais alta \cite{durant-smith} \cite{naive-forty}.

%Este tipo de classificador parte da ideia de que as informações presentes em um documento, utilizadas na determinação de sua classe, são independentes entre si. No caso de um documento de texto, assume-se que a presença ou ausência de um termo - uma palavra ou uma sequência de palavras - é independente da presença ou ausência de qualquer outro. Definida esta hipótese, a probabilidade de que um documento \emph{d} seja de uma classe \emph{c} é tal que

%\begin{equation}
%\label{eq1}
%\ensuremath{P(c|d) \propto P(c)\prod_{k=1}^{t_d}P(t_k|c)}  
%\end{equation} 

%onde \ensuremath{P(t_k|c)} é a probabilidade condicional do termo \ensuremath{t_k}  ocorrer em um documento da classe \ensuremath{c},  \ensuremath{P(c)} é a probabilidade a priori de um documento qualquer pertencer à classe \ensuremath{c} e \ensuremath{t_d} é o número de termos em \emph{d} \cite{stanford-IRbook}.

%As probabilidades envolvidas na relação \ref{eq1} contêm integrais díficeis ou mesmo impossíveis de se calcular analiticamente. Para calcular \ensuremath{P(c|d)}, portanto, utiliza-se aproximações obtidas através de técnicas de amostragem. Uma destas técnicas, comum na literatura de Aprendizado de Máquina e empregada neste projeto, é a amostragem de Gibbs. Em uma iteração da amostragem, a técnica condiciona as probabilidades calculadas para um documento \ensuremath{k} às classificações obtidas para os \ensuremath{k - 1} documentos anteriores \cite{resnik-gibbs}. A técnica está descrita, de forma básica, no algoritmo \textbf{X (o algorithm2e.sty tá dando pau)}.
                             % The basic idea in Gibbs sampling is that, rather than probabilistically picking
%the next state all at once, you make a separate probabilistic choice for each of the k dimensions, where each
%choice depends on the other k − 1 dimensions.


%A hipótese de independência entre os termos, razão pela qual o Naïve Bayes tem este nome\footnote{Naïve é uma palavra de origem francesa que significa "ingênua"}, simplifica bastante a estrutura da informação contida nos documentos. Ainda assim, o classificador costuma apresentar boas performances em categorização de textos, sendo utilizado, por exemplo, como base metodológica para alguns filtros de \emph{spam} \cite{paul}. Para melhorar o desempenho do Naïve Bayes, é comum fixar um conjunto de documentos previamente classificados de forma correta e utilizar a informação sobre suas classes na determinação das classes de outros documentos. Ao conjunto de documentos previamente classificados, dá-se o nome de \textbf{conjunto de treinamento}; ao conjunto de documentos a serem classificados, \textbf{conjunto de teste}.

 %Através do uso de Amostragem de Gibbs, a cada iteração se obtém uma aproximação melhor O número de amostras coletadas envolvendo cada classe foi fixado em 500. 
 
\subsection{SVMs}
\label{subsection:SVMs}

SVMs são uma família de métodos que utilizam uma abordagem geométrica para classificação. Eles são fundamentalmente utilizados em problemas de classificação envolvendo duas classes, mas podem ser adaptados para problemas mais complexos. Nesta seção, serão apresentados apenas os princípios de funcionamento de SVMs para duas classes, mais comuns na literatura. Para um aprofundamento sobre SVMs aplicados a problemas com mais de duas classes, recomenda-se a leitura do livro de Aprendizado de Máquina de Christopher Bishop \cite{bishop}.

Dado um conjunto de \ensuremath{n} pontos \ensuremath{\{x_i, y_i\}}, onde \ensuremath{x_i} é a representação vetorial de um documento \ensuremath{d} em um espaço euclidiano \ensuremath{\mathbb{R}^M} e \ensuremath{y_i} é sua respectiva classe, \ensuremath{y_i \in \{-1, 1\}}, um SVM deve decidir a classe \ensuremath{y} de um novo documento representado pelo vetor \ensuremath{x}. Para isso, assume-se que há pelo menos um hiperplano \ensuremath{\theta_0} que separa os pontos  com \ensuremath{y_i = 1} daqueles com \ensuremath{y_i = -1}. Um hiperplano \ensuremath{\theta_0} pode ser definido como o conjunto de pontos \textbf{x} que satisfazem

\begin{equation}
\label{function:svm}
\textbf{x} \ensuremath{\cdot} \textbf{w} + \ensuremath{b} = 0
\end{equation}

\textbf{w} é a normal ao hiperplano e \ensuremath{|b|||}\textbf{w}\ensuremath{||} é sua distância perpendicular à origem \cite{mono-puc}. A ideia é escolher os parâmetros \textbf{w} e \ensuremath{b} que maximizem a soma das distâncias dos hiperplanos \ensuremath{\theta_1} (vide equação \ref{theta1:svm}) e \ensuremath{\theta_{-1}} (vide equação \ref{thetamenos1:svm}) ao hiperplano \ensuremath{\theta_0}. 

\begin{equation}
\label{theta1:svm}
\textbf{x} \ensuremath{\cdot} \textbf{w} + \ensuremath{b} = 1
\end{equation}

\begin{equation}
\label{thetamenos1:svm}
\textbf{x} \ensuremath{\cdot} \textbf{w} + \ensuremath{b} = -1
\end{equation}

\ensuremath{\theta_1} e \ensuremath{\theta_{-1}} podem ser encontrados minimizando-se 

\begin{equation}
\label{optim:svm}
\ensuremath{\frac{1}{2}||}\textbf{w}\ensuremath{||^2}
\end{equation}

Para realizar esta otimização mais facilmente, o problema pode ser remodelado com multiplicadores de Lagrange \ensuremath{\{\alpha_i\}}, \ensuremath{1 \leq i \leq n}, levando à Equação \ref{lagrange} \cite{mono-puc}. Busca-se, então, a minimização desta equação com relação a \textbf{w} e \ensuremath{b} e maximização com relação a \ensuremath{\{\alpha_i\}}, com todo \ensuremath{\alpha_i \geq 0}.

\begin{equation}
\label{lagrange}
\ensuremath{L(\alpha, b,}\textbf{w}\ensuremath{) = \frac{1}{2} ||}\textbf{w}\ensuremath{||^2 - \sum_{i = 1}^n \alpha_i\big[y_i(x_i \cdot}\textbf{w}\ensuremath{+ b) -1 \big]} %x_i \cdot} \textbf{w} + \ensuremath{b} = 1 para \ensuremath{y_i = 1}
\end{equation}

Após a obtenção dos valores de \ensuremath{\{\alpha_i\}} que maximizam \ref{lagrange}, a obtenção da classe \ensuremath{y} de um documento representado por um vetor \ensuremath{x} é dada pelo sinal do somatório

\begin{equation}
\label{result:svm}
\ensuremath{y(x) = sinal\bigg(\sum_{i = 1}^n \alpha_iy_i(x_i \cdot x) + b\bigg)} %x_i \cdot} \textbf{w} + \ensuremath{b} = 1 para \ensuremath{y_i = 1}
\end{equation}

Esta solução funciona em casos nos quais os pontos \ensuremath{\{x_i, y_i\}} são linearmente separáveis - ou seja, obedecem à restrição

\begin{equation}
\label{restr2:svm}
\ensuremath{y_i(x_i \cdot} \textbf{w} + \ensuremath{b -1) \geq 0,\ i = 1,...,n}
\end{equation}

Quando esses pontos não são linearmente separáveis, essa metodologia precisa ser ajustada, modelando a classificação errônea de documentos. Isto envolve a introdução de \ensuremath{n} variáveis frouxas \ensuremath{\epsilon_i}\footnote{Do inglês \emph{slack variables}.}, uma para cada ponto \ensuremath{(x_i, y_i)}. \ensuremath{\epsilon_i = 0} se \ensuremath{y(x_i) = y_i} e \ensuremath{\epsilon_i = |y_i - y(x_i)|} em caso contrário \cite{bishop}. O SVM deve, neste caso, minimizar

\begin{equation}
\label{nonlin:svm}
\ensuremath{C\sum_{i=1}^n\epsilon_i + \frac{1}{2}||}\textbf{w}\ensuremath{||^2}
\end{equation}

em que \ensuremath{C} é um parâmetro responsável por controlar o compromisso entre a penalidade das variáveis frouxas e a distância máxima dos hiperplanos \ensuremath{\theta_1} e \ensuremath{\theta_{-1}} ao hiperplano \ensuremath{\theta_0}. Na modelagem com multiplicadores de Lagrange, a Equação \ref{lagrange} deve ser otimizada de tal forma que todo \ensuremath{\alpha_i} deve ser maximizado obedecendo à restrição \ensuremath{0 \leq \alpha_i \leq C}. Desta forma, também se obtém a Equação \ref{result:svm} para determinação da classe de um novo documento.

SVMs não foram utilizados em nenhum experimento deste projeto, mas fazem parte da metodologia de alguns dos trabalhos revisados.
 
%     We shall assume for the moment that the training data set is linearly separable in
%feature space, so that by definition there exists at least one choice of the parameters
%w and b such that a function of the form (7.1) satisfies y(xn ) > 0 for points having
%tn = +1 and y(xn ) < 0 for points having tn = −1, so that tn y(xn ) > 0 for all
%training data points.


%Suppose some given data points each belong to one of two classes, and the goal is to decide which class a new data point will be in. In the case of support vector machines, a data point is viewed as a p-dimensional vector (a list of p numbers), and we want to know whether we can separate such points with a p − 1-dimensional hyperplane. This is called a linear classifier. There are many hyperplanes that might classify the data. One reasonable choice as the best hyperplane is the one that represents the largest separation, or margin, between the two classes. So we choose the hyperplane so that the distance from it to the nearest data point on each side is maximized. If such a hyperplane exists, it is known as the maximum-margin hyperplane and the linear classifier it defines is known as a maximum margin classifier.

% cada documento do corpus é representado como um vetor \ensuremath{x} em um espaço euclidiano \ensuremath{\mathbb{R}^M}. Na etapa de treinamento, \ensuremath{N} documentos, com suas respectivas classes \ensuremath{c_0} ou \ensuremath{c_1}, são apresentadas ao classificador, que deve aprender uma função e a classificação envolve o aprendizado de uma função que corte este hiperplano, dividindo os pontos em duas classes distintas. 





%\chapter{Principais trabalhos e \emph{datasets} estudados}

%\section{Bitterlemons}
%\label{bitter}

%Falar de tudo e dos pre-processamentos

\chapter{Classificação por perspectiva: revisão e discussão}
\label{chap3}

A classificação de documentos de acordo com suas perspectivas é um tópico de pesquisa relativamente novo. Ele foi proposto como subproblema da área de Mineração de Opinião em 2008, pelas pesquisadoras Bo Pang e Lillian Lee \cite{omsa}. De acordo com elas, esse subproblema se diferencia da classificação por opinião por não configurar as classes como pólos (\emph{positivo}, \emph{negativo} etc.). De fato, a classificação por perspectiva visa a separar os documentos de um corpus de acordo com pontos de vista diferentes, como \emph{pró-Israel} ou \emph{pró-Palestina}. Embora, em certo nível, isso também polarize o corpus, já que normalmente os pontos de vista são antagônicos, dividi-lo de acordo com suas perspectivas é uma tarefa mais subjetiva do que aquela envolvendo opiniões. Um bom exemplo da diferença entre um ponto de vista e uma opinião, no contexto dos trabalhos revisados e de acordo com a \emph{survey} de Pang e Lee, é o seguinte: \emph{'Matrix' é um filme excelente} (opinião sobre um objeto específico, mais pontual) e \emph{a paz mundial dificilmente será alcançada} (ponto de vista sobre um tema; algo que pode implicar em uma série de opiniões alinhadas com este pensamento). Apesar de ter sido proposto formalmente apenas em 2008, uma conferência importante para a área de Mineração de Opinião, a \emph{AAAI\footnote{\emph{Association for the Advancement of Artificial Intelligence}.} Conference}, lançou, na edição de 2010, uma trilha específica para trabalhos nesta direção: \emph{Sentiment and Perspective} \cite{aaai2010}. Isso indica que o subproblema deve se fortalecer como área de pesquisa nos próximos anos.

A \emph{survey} de Pang e Lee, uma das principais referências para este projeto, elenca três trabalhos que fazem classificação por perspectiva, publicados entre os anos de 2004 e 2006. Assim como nesse \emph{survey}, esta monografia enfoca trabalhos recentes, publicados entre 2000 e 2010. A metodologia aplicada na busca por trabalhos a serem revisados pode ser resumida nos seguintes passos:

\begin{enumerate}
\item Consideração dos três trabalhos indicados na \emph{survey} de Pang e Lee;
\item Busca, no \emph{Google Scholar}, por trabalhos citados nesses artigos;
\item Busca, no \emph{Google Scholar}, por trabalhos que citam os artigos coletados nos itens anteriores;
\item Busca, na \emph{ACL Anthology}\footnote{A \emph{Association of Computational Linguistics}, \emph{ACL}, mantém o maior arquivo digital envolvendo trabalhos de Linguística Computacional, o que inclui Mineração de Opinião (http://aclweb.org/anthology-new/).} pelas palavras \emph{perspective}, \emph{viewpoint}, \emph{politics}, \emph{political}, \emph{stance}, \emph{ideology} e \emph{ideological}. Estas palavras foram escolhidas por (1) funcionarem como sinônimos nos artigos coletados anteriormente ou (2) pela relevância em suas temáticas, escritos por especialistas;
\item Busca nos \emph{sites} dos eventos \emph{EMNLP\footnote{\emph{Empirical Methods in Natural Language Processing Conference.}}, NAACL\footnote{\emph{North American Chapter of the Association for Computational Linguistics Conference.}}, AAAI\footnote{\emph{Association for the Advancement of Artificial Intelligence Conference.}}} e \emph{CoNNL\footnote{\emph{Computational Natural Language Learning Conference.}}}, realizados entre 2000 e 2010, pelas mesmas palavras do item anterior. Esses eventos foram selecionados pela relevância na área de Mineração de Opinião. 
\end{enumerate}

Nem todos os trabalhos coletados, de acordo com essa metodologia, tinham necessariamente a ver com classificação por perspectiva. Alguns, como a tese de Alice Oh, se propõem modelar computacionalmente o conceito de perspectiva \cite{alice-oh}; outros, como o artigo de Laver et al., se propõem a criar uma escala ideológica e comparar documentos de acordo com ela \cite{wordscores}.  Por este motivo, foi feita uma filtragem que resultou em 14 trabalhos a serem revisados. Na seção \ref{freqs:revisao}, essa revisão é apresentada. Quase todos os trabalhos consideram, ainda que não exclusivamente, as contagens de palavras dos documentos na classificação. Além disso, eles divergem bastante quanto à escolha de outras características dos documentos. Por estes motivos, a seção \ref{freqs:experim} discute a relação entre a classificação baseada em contagens de palavras e o emprego delas por diferentes perspectivas. Por fim, na seção \ref{freqs:conc}, apresentam-se conclusões sobre os conteúdos expostos nas seções \ref{freqs:revisao} e \ref{freqs:experim}. 

\section{Trabalhos Revisados}
\label{freqs:revisao}

Todos os trabalhos revisados nesta seção apresentam uma metodologia em comum: eles organizam um corpus de documentos, definem em que perspectivas ele se divide, decidem como os documentos serão representados e pré-processados e, em seguida, classificam-nos utilizando alguma técnica - na maioria dos casos, SVMs ou Naïve Bayes. Por questões de escopo, portanto, esses serão os pontos enfocados nesta seção. Inicialmente, os principais pontos de cada trabalho serão apresentados em subseções específicas. Por fim, na subseção \ref{compara}, será apresentada uma comparação entre eles, considerando os pontos enfocados por esta revisão. %Os trabalhos serão apresentados de maneira comparativa, de acordo com estes pontos, a fim de evidenciar as diferentes decisões de projeto que cada um toma. De todo modo, para prevenir uma visão muito fragmentada de cada um deles, foi construída a Tabela \ref{freqs:todos}, na qual cada trabalho pode ser visto, de forma  

\subsection{Posicionamentos sobre o conflito Israel-Palestina}

O trabalho de Lin et al. classifica artigos, escritos por especialistas, sobre o conflito Israel-Palestina \cite{lin-et-al2006}. Os artigos foram extraídos do \emph{site} Bitterlemons\footnote{http://www.bitterlemons.org/} e estão divididos entre \emph{Palestinian} ou \emph{Israeli}. Os primeiros defendem pontos de vista pró-Palestina; os segundos, pró-Israel. Essa divisão é explorada mais à frente, de modo que o objetivo desse trabalho é classificar os documentos como pró-Palestina ou pró-Israel.

Inicialmente, os documentos são representados como listas de palavras reduzidas a seus radicais. Isso significa, por exemplo, que as palavras \emph{political} e \emph{politics} são representadas através de um mesmo termo, pois possuem o mesmo radical \emph{politic}. Os classificadores Naïve Bayes e um SVM são então aplicados, explorando as contagens desses radicais em cada documento. A taxa de acerto obtida com o SVM variou entre 81.48\% e 97.24\%; com o Naïve Bayes, ela variou entre 84.85\% e 99.09\%. Essas variações advêm de diferentes divisões entre os conjuntos de treinamento e teste.

Por fim, o trabalho de Lin et al. propõe um novo classificador, denominado \emph{Latent Sentence Perspective Model}, ou simplesmente \textbf{LSPM}. Ele consiste em uma variação do Naïve Bayes que, em vez de considerar documentos como listas de palavras, os representa como listas de senteças. A hipótese assumida por ele é a seguinte: nem todas as sentenças carregam um ponto de vista, de modo que o classificador, antes de definir a classe de um documento, deve selecionar quais delas carregam palavras relevantes. As sentenças com termos mais comuns, considerando o corpus como um todo, não contribuem diretamente para a classificação. A taxa de acerto obtida com o LSPM variou entre  86.99\% e 94.93\%, também como consequência de diferentes divisões entre os conjuntos de treinamento e teste. Para essas divisões, as taxas de acerto obtidas com o Naïve Bayes foram de 84.85\% e 93.46\%. O SVM não foi comparado diretamente com o LSPM. Lin et al. concluem, diante dessas comparações, que o LSPM apresenta desempenho superior ao do Naïve Bayes. Nenhum outro trabalho revisado, entretanto, faz uso desse classificador. O tutorial de Resnik e Hardisty sobre Naïve Bayes e LSPM sugere algumas equações para a implementação deste último \cite{resnik}, mas o próprio Resnik afirmou, em \emph{e-mail} endereçado à autora desta monografia, nunca ter conseguido replicar os resultados apresentados por Lin et al. \cite{email}. Aparentemente, portanto, o LSPM não é tão simples de implementar quanto o Naïve Bayes.
 
\subsection{Posicionamentos em um fórum de discussão política}

Mullen e Malouf propõem dois trabalhos que analisam um conjunto de \emph{posts} do fórum de discussão política Politics\footnote{http://politics.com}, escrito por cidadãos comuns dos Estados Unidos \cite{aaai-politics} \cite{malouf-taking_sides}. Em ambos, apesar de admitirem a diversidade de posicionamentos contidos no fórum, os autores dividem os documentos em apenas duas perspectivas, por uma questão de simplicidade: liberal e conservadora. Os documentos devem ser classificados, portanto, como alinhados com uma dessas duas ideologias.

No primeiro trabalho \cite{aaai-politics}, eles tratam cada documento como a união de todos os \emph{posts} de um mesmo usuário. Cada documento é representado como uma lista de palavras, e um Naïve Bayes é aplicado considerando suas contagens em cada documento. A taxa de acerto obtida, via validação cruzada de dez dobras, foi de 60.37\%. O tamanho do \emph{dataset} (apenas 185 documentos) e o uso parecido de palavras por ambas as ideologias foram apontados como alguns dos principais motivos para a obtenção desse resultado. Mullen e Malouf também sugerem que a presença de documentos menores, correspondentes a usuários do fórum que raramente postam, pode contribuir negativamente para o desempenho do Naïve Bayes. A ideia é que, como esses usuários participam muito pouco do fórum, as palavras em seus \emph{posts} não são suficientes para se consolidar uma perspectiva, tornando-os mais difíceis de se classificar. Restringindo a classificação a usuários que postaram no fórum pelo menos 20 vezes, a taxa de acerto obtida com o Naïve Bayes foi ligeiramente superior (61.38\%). Uma última observação desse trabalho indica o que pode estar acontecendo: usuários liberais citam falas de usuários conservadores em 62.2\% de seus \emph{posts}; analogamente, conservadores citam liberais em 77.5\%  de seus \emph{posts}. Embora o artigo não indique a proporção média entre essas citações e o restante dos documentos, é possível que elas estejam interferindo negativamente no aprendizado das perspectivas. A presença das perspectivas liberal e conservadora em um mesmo documento, com uma correspondendo às intenções do usuário e outra sendo citada, pode homogeneizar o uso de palavras no corpus, comprometendo a viabilidade da classificação.

Essa forma de interação entre os participantes do fórum é explorada pelo segundo trabalho de Mullen e Malouf \cite{malouf-taking_sides}, com o intuito de melhorar a classificação. Para isto, cria-se um grafo de co-citação em que cada vértice representa um participante e cada citação em uma fala a outra é indicada por uma aresta entre seus autores. Os participantes são agrupados de acordo com seus padrões de citação e, em seguida, as falas de cada grupo obtido são tratadas como um único documento. Aplica-se um Naïve Bayes a esta nova coleção de documentos, também explorando apenas suas contagens de palavras, e os resultados obtidos são propagados para todos os participantes de cada grupo. Essa metodologia atinge resultados significativamente melhores que aqueles obtidos no trabalho anterior: para participantes com mais de 500 palavras de fala no fórum, a taxa de acerto relatada é de 73\%. É válido salientar que, muito provavelmente, o sucesso dessa técnica é consequência do fato de que usuários alinhados ideologicamente não se citam muito. Se, além de citar seus oponentes, eles também se citassem, os padrões de citação seriam sempre muito parecidos, inviabilizando os agrupamentos adequados no grafo.

\subsection{Posicionamentos em \emph{blogs} políticos}

O trabalho de Jiang e Argamon constrói um corpus de documentos extraídos de \emph{blogs} sobre política escritos em inglês \cite{jiang-argamon}. Os \emph{blogs} foram escolhidos de acordo com cinco catálogos \emph{online}: BlogCatalog, Blogorama, eTalkingHead, Campaigns and Elections e Best of the Web Blogs\footnote{http://blogcatalog.com, http://blogarama.com, http://etalkinghead.com, http://campaignsandelections.com e http://blogs.botw.org.}. Cada \emph{blog} é associado à perspectiva liberal ou conservadora, divisão explorada mais à frente na classificação. Os documentos que constituem o corpus são apenas as páginas iniciais de cada um deles. A finalidade desse trabalho, portanto, é determinar o posicionamento político desses \emph{blogs}, considerando seu conteúdo mais imediato - suas páginas iniciais.

Inicialmente, Jiang e Argamon aplicam um SVM às páginas, considerando apenas a presença/ausência de palavras nos documentos. A taxa de acerto obtida foi de 81.92\% via validação cruzada de dez dobras. Em seguida, as páginas foram reduzidas a suas sentenças subjetivas, através de uma seleção baseada no dicionário semântico \emph{General Inquirer}\footnote{http://www.wjh.harvard.edu/~inquirer/}. Para uma sentença ser escolhida, ela deveria conter pelo menos duas palavras categorizadas, segundo o dicionário, como pertencentes às categorias \emph{Dor}, \emph{Hostilidade}, \emph{Prazer} ou \emph{Virtude}, dentre outras. A taxa de acerto, neste cenário, subiu para 83.28\%. 

Esse trabalho, por fim, busca extrair as expressões opinativas que mais se relacionam aos pontos de vista liberal e conservador, baseando-se nas sentenças subjetivas e em consultas ao dicionário \emph{General Inquirer}. O SVM é aplicado considerando a presença/ausência de palavras e, adicionalmente, a presença/ausência das expressões opinativas separadas para cada lado. A taxa de acerto neste cenário foi de 84.96\%. Apesar de melhorarem a taxa de acerto obtida inicialmente, essas metodologias são completamente dependentes da língua em que o corpus está escrito. Elas não são facilmente adaptáveis a línguas que não dispõem de dicionários como o \emph{General Inquirer} \emph{online}. %As metodologias empregadas são completamente dependentes da linguagem do corpus, por utilizarem dicionários desenvolvidos para o inglês. % português.


\subsection{Posicionamentos sobre o Presidente Bush e a Guerra do Iraque}


\subsection{Posicionamentos sobre a Pena de Morte}
%tese e artigo de greene

\subsection{Posicionamentos de congressistas sobre a legislação dos Estados Unidos}
%convote e disagree

\subsection{Posicionamentos em \emph{sites} de música e política}

\subsection{Posicionamentos em debates \emph{online}}

\subsection{Análise de discursos políticos e suas reais ideologias}

\subsection{Análise da relação entre a defesa de uma perspectiva e a escolha de palavras}


\subsection{Comparações}
\label{compara}
  
\section{Experimentos com L-LDA e Naïve Bayes}
\label{freqs:experim}

Se um classificador utiliza apenas as contagens de palavras dos documentos para identificar suas perspectivas, sua taxa de acerto é tão mais baixa quanto menos essas contagens mudam de uma perspectiva para outra. Apesar dessa relação ser evidente, não se conhece nenhum método para quantificá-la. Como o seu entendimento amplia a compreensão dos resultados obtidos com esses classificadores, esta seção se detém a ilustrá-la através de alguns experimentos. A ideia é comparar a forma como as palavras são usadas em dois \emph{datasets}: no primeiro, a taxa de acerto obtida com um classificador \textbf{Naïve Bayes}, considerando apenas contagens de palavras, deve ser alta; no segundo, baixa. Para a análise do uso das palavras, será utilizado o modelo de tópicos \textbf{L-LDA}. A escolha do Naïve Bayes advém do fato de que os \emph{datasets} explorados nesta seção não são muito grandes, e o desempenho desse classificador, nesses casos, tende a ser superior ao obtido com SVMs \cite{ng-jordan}. O uso do LSPM não foi cogitado, devido aos pontos discutidos na seção \ref{sec:aaai}.

%APRESENTE OS DOIS DATASETS. JUSTIFIQUE PQ REFEZ EXPERIMENTOS OU PQ N ESCOLHEU O AAAI E SIM OUTRO. - OK 

O primeiro \emph{dataset} estudado é o mesmo discutido na seção \ref{sec:lin-et-al}, composto de artigos sobre o conflito Israel-Palestina. A taxa de acerto obtida com um Naïve Bayes aplicado a esse corpus foi alta, variando entre 84.85\% a 93.46\%. Inicialmente, pensou-se em estudar também o corpus discutido na seção \ref{sec:aaai}, composto de \emph{posts} do fórum Politics.com. O Naïve Bayes não classificou muito bem seus documentos,  atingindo taxas de acerto entre 60.37\% e 64.48\%. Infelizmente, não foi possível obtê-lo mediante solicitação aos autores do artigo. Por este motivo, o segundo \emph{dataset} estudado provém de outro trabalho: o artigo de Thomas, Pang e Lee sobre classificação de perspectiva em debates políticos dos Estados Unidos \cite{get-out-the-vote}. Esse trabalho, cujo corpus está disponível na página de Lee\footnote{http://www.cs.cornell.edu/home/llee/data/convote.html}, é um dos três que não utilizam apenas contagens de palavras para a classificação, sendo detalhado no Capítulo \ref{outros}. Por ora, é suficiente informar que ele é composto de \textbf{8126} trechos de discursos em debates da \emph{House of Representatives}, um dos dois órgãos principais do poder legislativo federal dos Estados Unidos. Os documentos estão divididos de acordo com duas orientações políticas antagônicas: a republicana  (4044) e a democrata (4046). Como no artigo original outras propriedades dos documentos foram consideradas, foi necessário aplicar o Naïve Bayes a esses documentos, considerando apenas as contagens de palavras dos mesmos. No trabalho de Thomas, Pang e Lee, os resultados apresentados são bons - no experimento efetuado com o Naïve Bayes para esse capítulo, entretanto, a taxa de acerto obtida, via validação cruzada de dez dobras, foi de 54.01\%\footnote{A precisão obtida foi de 54.47\% e a métrica F1 foi de 51.42\%.}.  

%MOSTRE COMO OS MODELOS DE TÓPICOS SÃO USADOS NOS DOIS LADOS - ok

Em ambos os \emph{datasets}, cada documento foi associado a dois tópicos: um genérico, idêntico para todos eles, e outro referente à sua perspectiva. No primeiro corpus, essas perspectivas são pró-Israel ou pró-Palestina; no segundo, republicana ou democrata. Há portanto, em cada corpus, três tópicos diferentes. O uso de um tópico genérico associado a todos os documentos ajuda a identificar palavras muito comuns nos \emph{datasets}, independentemente de perspectiva. Essa é a diferença fundamental entre essa aplicação do L-LDA\footnote{A implementação de L-LDA utilizada está disponível no repositório online de Alexandre Passos (http://github.com/alextp).} e a simples contagem de palavras em documentos, dividida entre duas perspectivas. Esse tipo de contagem não evidencia que palavras são mais escolhidas em documentos escritos sob uma certa perspectiva e quais são muito utilizadas por todos eles - informação que colabora para um maior entendimento das taxas de acerto supracitadas, obtidas com um Naïve Bayes. 

%MOSTRE AS PALAVRAS ELENCADAS E DISCUTA

As dez palavras mais frequentemente associadas a cada tópico, retirando-se artigos, conjunções, preposições, advérbios e pronomes pessoais, estão listadas nas Tabelas 4.1 e \ref{freqs:tab2}. Apesar de pequenas, essas listagens sugerem interpretações sobre como as palavras de cada corpus são exploradas por suas diferentes perspectivas. Essas interpretações, essencialmente  subjetivas, ajudam a entender o comportamento do classificador Naïve Bayes aplicado aos dois \emph{datasets}.%colaboram parsubjetiva de se ampliar a discussão sobre os resultados obtidos com o Naïve Bayes.

%provêem informações subjetivas sobre a linguagem empregada nos corpora.
%As palavras extraídas a partir da aplicação de um L-LDA  Ainda assim, essas informações 

%Para ilustrar a relação que existe entre as taxas de acerto do Naïve Bayes e o uso de palavras nesses corpora, a

%As análises feitas a partir dessas pequenas listagens são essencialmente subjetivas, sugerindo interpretações %sugerindo  Essa análise não pretende  % A partir disso, espera-se compreender melhor o 

\begin{table}[h]
\label{tab1}
\centering
\begin{tabular}{| l | p{10cm} | }
\hline
\textbf{Tópico} & \textbf{Palavras} \\ \hline
\textbf{Genérico} & israel, palestinian, israeli, palestinians, state, one, two, israelis, political, right \\ \hline
\textbf{Pró-Israel} & sharon, palestinian, arafat, peace, israeli, prime, bush, minister, american, process \\ \hline
\textbf{Pró-Palestina} & palestinian, israeli, sharon, peace, occupation, international, political, united, people, violence \\ \hline
\end{tabular}
\caption{As dez palavras mais frequentemente associadas aos tópicos pró-Israel, pró-Palestina e genérico, de acordo com um L-LDA.}
\end{table}


Considerando o primeiro corpus, as palavras associadas às perspectivas pró-Israel e pró-Palestina remetem de imediato ao conflito travado entre essas duas nações. Parte delas, como \emph{palestinian} e \emph{israeli}, são bastante mencionadas em ambas as perspectivas, ainda que com propósitos diferentes: %israelenses, 

\begin{quote}

\emph{"The recent \textbf{Israeli} government decision to begin building extensive walls
around \textbf{Palestinian} is just one more example of how \textbf{Israeli} Prime
Minister Ariel Sharon is unable to deal with \textbf{Israeli} problems save
through his narrow security vision."}

{\small Retirado de "Peace in peaces", de Ghassan Khatib (pró-Palestina) - 10/06/2002}
\end{quote}

\begin{quote}

\emph{"The first conclusion that the Israeli political and security
establishment should learn and internalize after 18 months of
\textbf{Palestinian} Intifada, concerns the intensity of \textbf{Palestinian} blind
terrorism and guerilla warfare against the State of Israel."}

{\small Retirado de "The lessons Israel should learn", de Meir Pa'il (pró-Israel) - 29/04/2002}
\end{quote}

%\emph{"Meanwhile, the Israeli army was still suffering in South Lebanon as a result of the dilemma that \textbf{Sharon} put them in."}

%{\small Retirado de "Sharon settles accounts", de Mamdouh Nofal (pró-Palestina) - 28/01/2002}

Outras palavras, como \emph{bush} e \emph{occupation}, foram mais enfatizadas, respectivamente, pelas perspectivas pró-Israel e pró-Palestina, constando em apenas uma listagem. Os trechos abaixo evidenciam a importância, no período em que os documentos desse corpus foram escritos (2001 - 2005), do Governo Bush para Israel e da criação de um Estado para a nação palestina:
 %\textbf{Falar das figuras e de como elas deixam isso tão claro.} 

%\textbf{FIGURAS} % Para compreender melhor o uso das palavras pelos diferentes pontos de vista do corpus, elas foram processadas pelo software wordle22 , resultando nas figuras 7.1, 7.2 e 7.3. O tamanho das palavras nas imagens corresponde ao quanto elas se associam a cada tópico23 .

%Falar das ênfases no próprio lado/lado oposto e mostrar tb as palavras genéricas.
%Os trechos abaixo, separados por perspectiva, ilustram dois contextos diferentes para o uso dessas palavras

%\begin{quote}
%\emph{"The recent \textbf{Israeli} government decision to begin building extensive walls
%around \textbf{Palestinian} is just one more example of how \textbf{Israeli} Prime
%Minister Ariel Sharon is unable to deal with \textbf{Israeli} problems save
%through his narrow security vision."} 

%{\small Retirado de \emph{"Peace in pieces"}, de Ghassan Khatib - 10/06/2002}
%\end{quote}

%\begin{quote}
%\emph{"The first conclusion that the Israeli political and security
%establishment should learn and internalize after 18 months of
%\textbf{Palestinian} Intifada, concerns the intensity of \textbf{Palestinian} blind
%terrorism and guerilla warfare against the State of Israel."} 

%{\small Retirado de }

%\end{quote}


\begin{quote}
\emph{"\textbf{Bush} and his advisers, who have been critical of Clinton's deep involvement
in a failed peace process ever since taking office, nevertheless
understood at the time that peace in the Middle East should be beyond
politics in America, and that the US could not permit itself to turn its
back on an Israeli leader who was determined to make peace."} 

{\small Retirado de "Barak was willing, and so were US Jews", de Yossi Alpher (pró-Israel) - 15/07/2002}
\end{quote}

\begin{quote}
\emph{"But just as we were close to a complete
package that would have ended the \textbf{occupation} and established a
Palestinian state, Barak permitted Ariel Sharon's provocative visit to
Al Aqsa mosque, and launched his "revenge" on Palestinians."} 

{\small Retirado de "Guarding our legitimacy", de Samir Abdullah (pró-Palestina) - 14/10/2002}
\end{quote}

Associadas ao tópico genérico, também estão as palavras \emph{palestinian} e \emph{israeli}, bastante exploradas pelas duas perspectivas. Isso reflete a importância que esses termos têm para o corpus como um todo. Palavras relacionadas mais genericamente ao conflito Israel-Palestina, como \emph{state} e \emph{political}, também aparecem na listagem. %, evidenciando indica que essas palavras         


                     % A imagem 7.3 dá certo destaque a Lula e José Serra, mas também enfatiza outros termos, como governo e brasil, relacionados mais genericamente ao tema geral dos artigos: a política brasileira.


%Essas análises, avaliadas à luz das altas taxas de acerto obtidas com o Naïve Bayes, sugerem que os autores desse primeiro corpus defendem seus pontos de vista veementemente. Como eles são antagônicos, isso se reflete em enfoques diferentes a determinados termos, %empregos diferentes de palavras. % tornando a forma como empregam as palavras  defendem suas perspectivas utilizando palavras de forma suficientemente diferentedestacando ideias diferentes.

\begin{table}[h]
\label{freqs:tab2}
\centering
\begin{tabular}{| l | p{10cm} | }
\hline
\textbf{Tópico} & \textbf{Palavras} \\ \hline
\textbf{Genérico} & mr., speaker, bill, all, time, people, today, gentleman, federal, support \\ \hline
\textbf{Democrata} & bill, security, legislation, states, chairman, country, act, billion, million, law \\ \hline
\textbf{Republicano} & act, chairman, security, states, bill, legislation, 11, support, 9, system \\ \hline
\end{tabular}
\caption{As dez palavras mais frequentemente associadas aos tópicos republicano, democrata e genérico, de acordo com um L-LDA.}
\end{table}

%Considerando o primeiro corpus, as palavras associadas às perspectivas pró-Israel e pró-Palestina remetem de imediato ao conflito travado entre essas duas nações. Parte delas, como \emph{palestinian} e \emph{israeli}, são bastante mencionadas em ambas as perspectivas, ainda que com \textbf{intensidades} diferentes. \textbf{Falar das figuras e de como elas deixam isso tão claro.} 

%Falar q blz, aparecem palavras diferentes e tb há proporções diferentes, mas oq dá pra notar no geral eh q os termos envolvem pouca polêmica, e são usados de forma parecida. Isso, e o fato da amostra ser pequena, sugere q os debatentes n defendem as perspectivas republicana e democrata de forma suficientemente clara pra q um Naïve Bayes classifique corretamente.

Considerando o segundo corpus, tem-se que parte das palavras associadas às perspectivas republicana e democrata, como \emph{bill}, \emph{legislation}, \emph{states} e \emph{act}, tem mais a ver com o processo legislativo em si do que com alguma delas. Isso sugere que os debatentes não focam seus discursos na defesa de ideias republicanas ou democratas - ou, se o fazem, não utilizam as palavras de forma suficientemente diferente,  para que um Naïve Bayes classifique-os corretamente. Reforçando a ideia de que esse corpus é pouco polêmico, dificultando a identificação de perspectivas diferentes, foi observado que algumas das palavras listadas na Tabela \ref{freqs:tab2} são, muitas vezes, empregadas com conotações semelhantes por democratas e republicanos. Um exemplo é a palavra \emph{security}, como pode ser observado nos trechos abaixo

\begin{quote}
\emph{"Mr. speaker, I wholeheartedly agree that if we want to cut down on illegal immigration, we must improve border \textbf{security}. Just 2 weeks ago, an astute crane operator at the port of Los Angeles discovered 32 Chinese stowaways in a container that had just been unloaded from a Panamanian freighter."} 

{\small Retirado de discurso de Jane Harman, democrata - 09/02/2005 }

\end{quote}

\begin{quote}
\emph{"The fence remains incomplete and is an opportunity for aliens to cross the border illegally. This incomplete fence allows border \textbf{security} gaps to remain open.  We must close these gaps because they remain a threat to our national \textbf{security}."}

{\small Retirado de discurso de John Boozman, republicano - 02/10/2005}

\end{quote}

Duas palavras listadas apenas para a perspectiva republicana, \emph{11} e \emph{9}, podem sugerir um maior enfoque no episódio 11 de Setembro e seus desdobramentos. A palavra \emph{billion}, listada apenas para a perspectiva democrata, é muitas vezes utilizada para discutir gastos públicos, o que pode indicar um destaque maior para esse assunto

\begin{quote}
\emph{We know that all but one of the \textbf{9}/\textbf{11} hijackers acquired some type of U.S. identification documents. In fact, the 19 hijackers had 63 driver 's licenses among them.}

{\small Retirado de discurso de John Sullivan, republicano - 09/02/2005} 
\end{quote}

\begin{quote}
\emph{The Pomeroy bill would cost the treasury \$72 \textbf{billion} over 10 years, compared with the \$290 \textbf{billion} price tag of a full repeal through 2015, according to the joint committee on taxation.}

{\small Retirado de discurso de Earl Pomeroy, democrata - 12/04/2005}
\end{quote}

As palavras associadas ao tópico genérico, assim como aquelas associadas às perspectivas republicana e democrata, têm mais a ver com o processo legislativo e a estrutura dos debates do que com posicionamentos políticos. Alguns exemplos são \emph{bill}, \emph{mr.}, \emph{speaker} e \emph{gentleman}. Isso reforça a ideia de que é difícil identificar as perspectivas dos documentos desse corpus. %Associadas ao tópico genérico, há palavras referentes ao processo legislativo e aos debates, como \emph{bill}, \emph{mr.}, \emph{speaker} e \emph{gentleman}. Essas são as palavras mais comuns em todo o corpus - e, como pode ser observado, elas se assemelham às palavras elencadas para as perspectivas. Não apenas porque algumas delas, como blablala, estão listadas para os três tópicos, mas especialmente porque grande parte delas não evidencia um viés republicano ou democrata.%  Essa característica também é observada nas palavras listadas para as duas perspectivas, o que sugere a dificuldade de se identificar as perspectivas republicana ou democrática nos documentos desse segundo corpus. %  estão as palavras \emph{palestinian} e \emph{israeli}, bastante exploradas pelas duas perspectivas. Isso reflete a importância que esses termos têm para o corpus como um todo. Palavras relacionadas mais genericamente ao conflito Israel-Palestina, como \emph{state} e \emph{political}, também aparecem na listagem. 

As Tabelas 4.1 e \ref{freqs:tab2} sugerem que os republicanos e democratas estudados se expressam de forma mais parecida do que os autores do primeiro corpus. Talvez isso advenha do fato de que seus discursos fazem parte de debates, o que pode concentrar as discussões em torno de um conjunto muito específico de termos. No primeiro corpus, os autores não escreveram seus artigos como resposta direta a outros, o que pode colaborar para que eles se expressem de forma mais autêntica e veemente, focando seus textos na apresentação de seus pontos de vista. A aplicação do L-LDA, associando cada tópico a uma perspectiva diferente, se mostrou válida, evidenciando uma certa homogeneização no uso de palavras no segundo corpus, quando comparado com o primeiro.
%As taxas de acerto obtidas com um Naïve Bayes nesse segundo corpus foram baixas. % Ainda que as perspectivas republicana e democrata enfatizem palavras diferentes, conforme pode ser observado na Tabela \ref{freqs:tab2}, as listagens sugerem que elas empregam mais palavras em comum, no decorrer dos debates, do que os artigos  

As palavras das Tabelas 4.1 e \ref{freqs:tab2} foram representadas graficamente com o apoio do \emph{software} wordle\footnote{http://wordle.net/}. Apesar de se associarem aos tópicos de formas distintas, o \emph{software} não foi sensível o suficiente para captar essas diferenças. Por este motivo, essas imagens não constam nessa monografia. O número de vezes que cada palavra se associou a cada tópico, entretanto, pode ser consultado no \textbf{ANEXO BLA}. % foi 



É válido ressaltar que, a depender do \emph{dataset}, outras questões podem colaborar para um mau desempenho na classificação. Um conjunto de documentos com poucos exemplares, ou contendo poucas palavras, é um cenário onde a aplicação do Naïve Bayes pode não funcionar bem. Entretanto, esse não parece ser o caso dos corpora analisados nessa seção. De todo modo, quando não se obtém uma boa taxa de acerto com um classificador baseado em contagens de palavras, a investigação subjetiva de como cada perspectiva faz uso delas pode ajudar na compreensão do fenômeno. No segundo corpus analisado nesta seção, por exemplo, o uso de palavras que pouco têm a ver com as perspectivas republicana e democrata, muitas vezes destacadas pelas duas perspectivas, reforça a ideia de que contagens de palavras não são suficientes para classificação nesse cenário. %Apesar de não se quantificar a relação entre o uso de palavras no segundo corpus  sua classificação,  %A depender das conclusões retiradas, pode-se pensar em outras metodologias para a classificação dos documentos de acordo com suas perspectivas.

% mais relacionadas ao processo legislativo em si do que a suas ideias. % é preciso ter em mente que a amostra é pequena, de modo % se assemelham bastante, em termos temáticos, às palavras relacionadas ao tópico genérico. Isso reforça a i

% alta frequência de palavras como essas, empregadas pelos dois lados do debate, indica um cenário pouco polêmico, com menos \emph{banner words} e divergências. A palavra \emph{security}, por exemplo, fortemente associada às duas perspectivas, é utilizada de forma similar por ambas, como ilustrado na Tabela 4.6.

% Parte das palavras, como \emph{bill}, \emph{legislation}, \emph{states} e \emph{act},
%da Tabela 4.2 indicam que o vocabulário do segundo \emph{dataset} não é suficiente para distinguir as perspectivas Republicana e Democrata.

 % d da temática. %A alta frequência de palavras associadas às perspectivas, bem como a presença de \emph{banner words} importantes, conFiguram um bom cenário para o uso de métodos baseados em frequências de palavras. O desempenho de um Naïve Bayes na classificação deste \emph{dataset} será discutido mais à frente, ainda nesta seção.

%As listas de palavras da Tabela 4.2 indicam que o vocabulário do segundo \emph{dataset} não é suficiente para distinguir as perspectivas Republicana e Democrata. Parte das palavras, como \emph{bill}, \emph{legislation}, \emph{states} e \emph{act}, estão mais associadas ao processo legislativo \emph{per se} do que a alguma das perspectivas contidas nos documentos. A alta frequência de palavras como essas, empregadas pelos dois lados do debate, indica um cenário pouco polêmico, com menos \emph{banner words} e divergências. A palavra \emph{security}, por exemplo, fortemente associada às duas perspectivas, é utilizada de forma similar por ambas, como ilustrado na Tabela 4.6. Métodos baseados em frequências de palavras funcionam tão melhor quanto mais distintos forem os vocabulários empregados por cada perspectiva. Por este motivo, é esperado que suas taxas de acerto em \emph{datasets} como este não sejam altas.

%MOSTRE IMAGENS E TEXTINHOS E DISCUTA
 % supracitadas.  %Como esse tipo de contagem não objetiva separar palavras muito comuns em todo o corpus daquelas mais específicas de cada perspectivaconsidera palavras de \emph{background}, a visualização de palavras mais específicas para cada perspectiva é prejudicada. 


%O primeiro \emph{dataset} estudado é o \textbf{Bitterlemons}\footnote{A descrição deste \emph{dataset} encontra-se na seção \textbf{XXX} desta monografia}, composto de artigos pró-Israel e pró-Palestina. Cada documento foi associado a um tópico referente à sua perspectiva e outro genérico, idêntico para todos eles. Um modelo de tópicos do tipo L-LDA foi aplicado aos documentos assim anotados, agrupando palavras genéricas em torno do tópico genérico, pró-Israel em torno do tópico pró-Israel e pró-Palestina em torno do tópico pró-Palestina. As dez palavras mais fortemente associadas a cada um dos tópicos, excluindo-se \emph{stop words}, estão listadas na Tabela \ref{1}.




%O segundo \emph{dataset} estudado é o \textbf{Convote-Menor}, composto de colocações em debates da \emph{House of Representatives}, um dos dois órgãos principais do poder legislativo federal dos Estados Unidos. Os documentos foram marcados como sendo de parlamentares Republicanos ou Democratas, e como representando um posicionamento a favor ou contra a lei em pauta. Para este experimento, apenas a divisão entre Republicanos e Democratas foi considerada. O L-LDA foi aplicado a este \emph{dataset} de forma análoga ao primeiro experimento, e as dez palavras mais fortemente associadas a cada um dos tópicos - Genérico, Republicano e Democrata - estão listadas na Tabela 4.2. \emph{Stop words} também foram excluídas desta listagem.



% o que faria desse corpus u. A comparação entre a forma como as palavras são empregadas nos dois corpora

%O desempenho obtido com um classificador \textbf{Naïve Bayes} será analisado através da investigação de quais palavras se associam mais frequentemente a cada perspectiva de dois \emph{datasets}. 




%Para a seleção dessas palavras, o modelo \textbf{L-LDA} será utilizado.  %O que pode ser feito  essa relação.  




%mas não foi possível obtê-lo.


%Cada documento foi associado a um tópico referente à sua perspectiva e outro genérico, idêntico para todos eles. Um modelo de tópicos do tipo L-LDA foi aplicado aos documentos assim anotados, agrupando palavras genéricas em torno do tópico genérico, pró-Israel em torno do tópico pró-Israel e pró-Palestina em torno do tópico pró-Palestina. As dez palavras mais fortemente associadas a cada um dos tópicos, excluindo-se \emph{stop words}, estão listadas na Tabela \ref{1}.




% e explorada em alguns dos artigos estudados para este projeto, como o de Mullen e Malouf revisado na seção \ref{sec:aaai}

%Se pessoas com pontos de vista diferentes sobre um determinado tema escolhem palavras de forma muito parecida para se expressar, as contagens de palavras de seus textos podem não ser suficientes para diferenciar seus pontos de vista.


 %Embora óbvia, não se conhece nenhum método para mensurar essa ideia - ou seja, medir o quanto essas escolhas interferem na identificação das perspectivas de documentos. 


%nesse tipo de classificação, 



%Esta seção pretende ilustrar as relações existentes entre a identificação de perspectivas baseada em contagens de palavras e a forma como elas são empregadas nos documentos.  Este princípio foi ob


%Diante deste cenário,  se expressam através 



%A forma como as pessoas escolhem palavras para defender suas perspectivas interfere em classificadores baseados em contagens de palavras.




%m comum, e em proporções parecidas, 



%seus textos possuirão contagens de palavras semelhantes


 %a respeito de um tema, 



 %desempenho de um classificador baseado em contagens de palavras se um classificador utiliza contagens de palavras para identificar a             xperimentos com um classificador Naïve Bayes e um modelo de tópicos L-LDA
%ão conduzidos, a fim de se ilustrar as relações existentes entre o desempenho da classifica-
%ção por perspectiva e o emprego de palavras nos documentos


%Se um método utiliza apenas a frequência das palavras dos documentos para identificar suas perspectivas, é natural que a taxa de acerto seja tão mais baixa quanto menos essas frequências mudam de uma perspectiva para outra. Nesta seção, serão descritos experimentos que evidenciam o vocabulário contido em dois \emph{datasets}, as taxas de acerto obtidas na classificação dos documentos com um Naïve Bayes e a relação entre estas informações.

%O primeiro \emph{dataset} estudado é o \textbf{Bitterlemons}\footnote{A descrição deste \emph{dataset} encontra-se na seção \textbf{XXX} desta monografia}, composto de artigos pró-Israel e pró-Palestina. Cada documento foi associado a um tópico referente à sua perspectiva e outro genérico, idêntico para todos eles. Um modelo de tópicos do tipo L-LDA foi aplicado aos documentos assim anotados, agrupando palavras genéricas em torno do tópico genérico, pró-Israel em torno do tópico pró-Israel e pró-Palestina em torno do tópico pró-Palestina. As dez palavras mais fortemente associadas a cada um dos tópicos, excluindo-se \emph{stop words}, estão listadas na Tabela \ref{1}.

%O uso de um tópico genérico ajuda a identificar palavras de \emph{background}, comuns no corpus independentemente de perspectiva. Esta é a diferença fundamental entre o uso de um L-LDA e a simples contagem de palavras em documentos pró-Israel e pró-Palestina. Como esse tipo de contagem não considera palavras de \emph{background}, a visualização de palavras mais específicas para cada perspectiva é prejudicada.





%As palavras \emph{11} e \emph{9}, associadas à perspectiva Republicana,   Tabela \ref{2} também indica uma particularidade Republicana: o enfoque no episódio 11 de Setembro como base para argumentação a favor de novas leis. os dois lados em relação a algumas questões, como segurança.  

%distingue suficientemente duas perspectivas


%segundo \emph{dataset} apresenta poucas palavras especificamente associadas às perspectivas Republicana e Democrata. Parte das palavras listadas, como \emph{bill} e \emph{act}, estão diretamente associadas ao processo legislativo, independentemente dos pontos de vista debatidos \textbf{ilustrar com um exemplo}. 


%\textbf{Catar exemplos com msms palavras para perspectivas diferentes e palavras diferentes ressaltando lados opostos.}

%\textbf{Os exemplos... blablabla.}

%As palavras extraídas a partir da aplicação de um L-LDA provêem informações subjetivas sobre a linguagem empregada nos corpora. Ainda assim, essas informações ajudam a entender o comportamento do classificador Naïve Bayes aplicado aos dois \emph{datasets}. Para o \textbf{Bitterlemons}, as taxas de acerto obtidas variaram entre 73.46\% e 98.98\%, a depender da divisão entre os conjuntos de treinamento e teste; para o \textbf{Convote-Menor}, entre 48.73\% e 54.17\%. Não é trivial quantificar a relação entre essas taxas de acerto e a linguagem dos corpora - mas, como o Naïve Bayes utiliza apenas a distribuição das palavras para inferir a perspectiva dos documentos, é evidente que a escolha do vocabulário contribui para a qualidade da classificação.

%alta proporção de artigos, conjunções e pronomes em detrimento de possíveis \emph{stigma words} e \emph{banner words}, representadas por verbos, adjetivos e substantivos. Como consequência, a taxa de acerto do Naïve Bayes, treinado com 80\% do corpus e testado nos 20\% restantes, foi de apenas 48.92\%. 


%Para todos os experimentos, o número de iterações do L-LDA foi fixado em 100; o número de iterações do Naïve Bayes, em 500. O primeiro \emph{dataset} é composto de 594 documentos: 297 pró-Israel e 297 pró-Palestina. Ele contém 465.422 palavras, 14.500 delas únicas. O segundo \emph{dataset} é composto de 983 documentos: 487 escritos por Democratas e 496 escritos por Republicanos. Ele contém 359.761 palavras, 13.025 delas únicas. A implementação de L-LDA utilizada nestes experimentos está disponível em \cite{top-llda}; a de Naïve Bayes, em \cite{alibezz-nb}.





 %Investigar o vocabulário de um corpus, , pode ser interessante para verificar se sua uniformidade, ainda que em parte, está relacionada à má classificação obtida.
 




%e associam mais fortemente  para cada perspectiva, bem como aquelas que co-ocorrem em lados distintos do corpus, uniformizando o vocabulário.
%Para mostrar que uma análise prévia do vocabulário do corpus não é suficiente para recomendar - ou desaconselhar - o uso de classificadores baseados em frequências de palavras, um experimento envolvendo um modelo de tópicos foi executado. Em vez de contar as frequências das palavras para cada perspectiva presente em um \emph{dataset}




 %disponível em \cite{top-llda} aplicada a três \emph{datasets}.\textbf{explicar a fonte do naive bayes, os resultados, o uso de stop words, as iterações.}% O primeiro, composto de artigos extraídos do site bitterlemons.org, foi classificado com um Naïve Bayes padrão no artigo \cite{lin-et-al2006}. As taxas de acerto obtidas, com o uso do Naïve Bayes, variaram entre 84.85\% e 93.46\% para os experimentos elencados nesse artigo. O segundo, \textbf{definir dataset, espero que o politics.com}, também foi classificado com um Naïve Bayes padrão em \cite{malouf-taking_sides} - mas as taxas de acerto foram bem mais baixas: \textbf{X\%}. Para o experimento com o L-LDA, todas as palavras contidas nos documentos, incluindo \emph{stop words} como \emph{the}, foram consideradas, resultando em uma análise das palavras independente do pré-processamento executado nos \emph{datasets} nos dois artigos citados.

%\textbf{Imagens, desempenhos.}

%\textbf{reafirme os resultados acima.} Diante disso, uma estratégia mais indicada é começar o processo de classificação do \emph{dataset} com um classificador baseado, inicialmente, apenas nas frequências das palavras do corpus. Se a taxa de acerto atingida for menor do que a desejada, explora-se outras características do \emph{dataset}. Nesta seção, serão descritas apenas as técnicas de classificação dos artigos que associaram, explicitamente, a uniformidade das palavras no corpus ao mau desempenho de classificadores desse tipo\footnote{Técnicas de classificação desenvolvidas em outros artigos serão apresentadas em outras seções.}.

%\section{Aumentando as taxas de acerto: técnicas empregadas na literatura}

%Dos artigos estudados para este projeto, três associaram o mau desempenho obtido com Naïve Bayes e SVM padrão à homogeneização do vocabulário contido no corpus: dois de Tony Mullen e Robert Malouf sobre citações entre documentos \cite{malouf-taking_sides} e \cite{aaai-politics} e um de Miles Efron envolvendo citações a \emph{sites}\cite{efron}. Não foi possível executar experimentos com o L-LDA em nenhum dos \emph{datasets} utilizados nesses trabalhos, pois eles não estão disponíveis na Web nem conseguiram ser obtidos mediante pedido, por \emph{e-mail}, aos autores dos artigos. Por este motivo, esta seção se limitará a descrever as técnicas utilizadas nesses trabalhos para melhorar as taxas de acerto na classificação dos corpora.

%Em dois artigos \cite{malouf-taking_sides} e \cite{aaai-politcs}, o \emph{dataset} estudado é o \textbf{Politics}: um conjunto de debates políticos dos Estados Unidos. Os dois trabalhos visavam a classificar discursos de 185 participantes da discussão de acordo com suas orientações políticas: Esquerda ou Direita. Cada participante era representado por um único documento, resultante da concatenação de todas as suas falas nos debates. Aplicando um Naïve Bayes na coleção de documentos, as taxas de acerto obtidas foram de 60.37\% \cite{aaai-politics} e 63.59\% \cite{malouf-taking_sides}. A análise do \emph{dataset} indica que 62.2\% das falas de Esquerda mencionam trechos de falas de Direita; quanto às falas de Direita, 77.5\% delas mencionam falas de Esquerda \cite{aaai-politics}.

%\begin{figure}[t]
%  \centering % este comando é usado para centralizar a Figura
%  \includegraphics[width=8cm, height=6cm]{taking-sides-graph.png}\\
%  \caption{Grafo extraído de \cite{malouf-taking_sides}. Os vértices são usuários e as arestas, citações entre eles. Arestas mais escuras indicam frequência de citação mais alta.}
%  \label{fig:malouf}
%\end{figure}


%Padrões de citação a \emph{sites} também foram investigados \cite{efron}. No artigo de Miles Efron, os experimentos envolvem os \emph{datasets} \textbf{Political-Orientation} e \textbf{Artists}. O primeiro é composto de artigos políticos Estadunidenses de Direita ou Esquerda e o segundo é composto de textos sobre artistas musicais, divididos entre as categorias Alternativo e Popular. As taxas de acerto relatadas para as aplicações de Naïve Bayes nestes corpora foram de, respectivamente, 64.71\% e 50.1\%. A fim de melhorar as taxas de acerto na classificação dos documentos, o artigo estima a perspectiva de cada um deles avaliando a probabilidade de que eles sejam co-citados, na Web, com documentos referência de cada perspectiva. Estes documentos de referência são agrupados em listas de URLs. Esta metodologia, aplicada a \textbf{Political-Orientation}, resultou em uma taxa de acerto de 94.1\%. Para \textbf{Artists}, a taxa de acerto máxima obtida foi de 88.84\%.

%Todos os artigos que propuseram metodologias para melhorar a classificação em corpora com vocabulários muito uniformes consideraram algum esquema de citação. Se o corpus for composto de documentos pouco citados na Web, o esquema do artigo de Miles Efron pode não trazer nenhuma melhoria significativa ao problema. Se os autores dos documentos não mencionarem significativamente um ao outro, como aconteceu nos estudos com o \emph{dataset} \textbf{Politics}, a metodologia explorada por Tony Mullen e Robert Malouf não é recomendada. Essa área de pesquisa, portanto, ainda tem alguns problemas em aberto - e, considerando que estes três artigos foram escritos há mais de quatro anos, ela não parece atrair muita atenção. A justificativa para isso, provavelmente, tem a ver com o fato de que uniformidade no vocabulário não é um problema tão comum em Mineração de Perspectiva. No \textbf{capítulo X}, técnicas envolvendo comparação entre documentos, que podem trazer benefícios à classificação de \emph{datasets} difíceis - e neste caso também se enquadram aqueles em que o vocabulário é mais uniforme -, serão discutidas.

\section{Conclusões}
\label{freqs:conc}

%INTERNACIONALIZÁVEEEEEL!!

%classif. com contagens rocks, l-lda rocks, qto mais diferente melhor, mas pelos resultados obtidos na revisão, parece q eh bem por aí.

A classificação baseada em contagens de palavras, ou em alguma das variações mencionadas na introdução desse capítulo, assume a seguinte hipótese linguística: a quantidade de vezes que uma palavra é mencionada em um documento está diretamente relacionada a seu enfoque \cite{teubert}. Como consequência, esse método funciona melhor em \emph{datasets} nos quais o emprego de palavras varia significativamente por perspectiva. Esse parece ser o caso da maioria dos artigos estudados para este capítulo: cinco de seis trabalhos apresentaram uma boa taxa de acerto, considerando apenas contagens de palavras. Outros três trabalhos, apresentados no capítulo \ref{outros}, também fazem uso dessa informação - mas a associam a outras propriedades dos documentos classificados. De todo modo, a contagem de palavras mostrou ser a característica mais \textbf{essencial} à classificação por perspectiva, reforçando a ideia de que essa hipótese linguística é válida.


Nesse capítulo, foram revisados dois trabalhos que utilizam apenas essas contagens para classificar seus \emph{datasets} de estudo - eles são, inclusive, os mais citados dentre os treze analisados nessa monografia.

O primeiro, de Lin et al., apresenta um novo modelo para classificação (o LSPM), destacando-se dos demais artigos por não enfocar em SVMs ou Naïve Bayes. Diferentemente desses classificadores, o LSPM considera que apenas uma parte das sentenças de cada documento realmente apresenta um ponto de vista, e gera palavras mais específicas para cada uma delas. As demais frases concentram palavras mais genéricas, que poderiam ser empregadas da mesma forma por todos os lados do corpus. Por esse motivo, elas não contribuem diretamente com a classificação. O modelo apresenta uma boa taxa de acerto, mas não é tão estudado - nem tão trivial de implementar - quanto um Naïve Bayes. O trabalho é muito citado por ser um dos primeiros a tentar classificar documentos de acordo com suas perspectivas, e o \emph{dataset} analisado (artigos pró-Israel e pró-Palestina) também foi estudado por outros pesquisadores\footnote{Verificar \textbf{ANEXO BLA}.}.

O segundo, de Mullen e Malouf, é muito citado por ser um dos poucos que trabalha com um \emph{dataset} tão informal (\emph{posts} em um fórum sobre política), apresentando discussões interessantes sobre as dificuldades envolvidas no processo de classificá-lo. As taxas de acerto obtidas na classificação foram baixas, e um dos motivos possíveis, de acordo com o exposto no trabalho, é a quantidade de citações a textos escritos sob uma perspectiva diferente daquela que o autor defende. Isso \emph{mistura} contagens relativas a perspectivas distintas em um mesmo documento, homogeneizando-os sob o ponto de vista do classificador. 

O artigo de Lin et al. também apresentou uma boa taxa de acerto utilizando um Naïve Bayes, diferentemente do estudo de Mullen e Malouf. Apesar do corpus analisado por esses últimos ser pequeno, os autores indicam que a forma como cada perspectiva se apropria das palavras interfere na qualidade da classificação.  A fim de ampliar a compreensão sobre essa interferência, foram conduzidos dois experimentos envolvendo um Naïve Bayes e um L-LDA, aplicados a \emph{datasets} para os quais a classificação funcionou de forma diferente. O primeiro corpus considerado foi aquele estudado por Lin et Al., para o qual as taxas de acerto obtidas com um Naïve Bayes foram altas; o segundo, dado que não foi possível ter acesso aos documentos analisados por Mullen e Malouf, foi um conjunto de trechos de discursos em debates da \emph{House of Representatives}, órgão legislativo dos Estados Unidos. A taxa de acerto obtida com um Naïve Bayes nesse último corpus foi bastante baixa, tornando-o ideal para o objetivo dos experimentos. 

Nos experimentos com o L-LDA, cada documento foi associado a dois tópicos: um genérico e outro correspondente à sua perspectiva. Mesmo sabendo que em um Naïve Bayes não há distinção entre palavras genéricas e específicas, optou-se por essa divisão de tópicos por conta do objetivo da aplicação do L-LDA: a visualização parcial de que termos são mais enfocados por cada perspectiva. Essa informação sugere que há uma certa homogeneização nos enfoques do segundo \emph{dataset}, quando comparado com o primeiro. Considerando que quão mais homogêneo é o emprego de palavras por perspectiva, pior é o desempenho dos classificadores baseados em contagens de palavras, a informação obtida com o L-LDA ajuda a compreender porque a taxa de acerto obtida para o segundo \emph{dataset} foi tão mais baixa que aquelas obtidas para o primeiro. É importante frisar que não se encontrou \textbf{nenhum outro trabalho} que faça uso de um L-LDA para compreender, ainda que parcialmente, como certos termos são enfocados por diferentes perspectivas. Apesar de outros fatores contribuírem para o mau desempenho de uma classificação, como um número muito pequeno de documentos no \emph{dataset}, a investigação do emprego de palavras, quando as taxas de acerto obtidas não são boas, amplia a compreensão do corpus analisado - o que pode ser útil no momento de se pensar em outras estratégias para melhorar a classificação.  

%é interessante, para um classificador baseado em contagens de palavras, que elas variem o quanto for possível de uma perspectiva para outra, evidenciando os diferentes enfoques e pontos de vista dos autores dos documentos, a informação obtida com o L-LDA explica, ainda que  
 %, para compreender melhor como os autores dos documentos se expressam. % Se as palavras são empregadas de modo muito parecido por todas as perspectivas, isto justifica, ainda que em parte, o mau desempenho obtido. 

% Considerando que a classificação dos dois \emph{datasets} baseou-se apenas em contagem de palavras, a informação obtida com o L-LDA ampliou a discussão %m isso, e considerando a taxa de acerto obtida com o Naïve Bayes em cada caso,  % facilitou a visualização sobre aquilo que é mais específico de cada perspectiva. Como o uso do L-LDA objetivava apenas a uma 

%de modo que % A ideia era ilustrar, através de um L-LDA, o emprego de palavras em um corpus bem classificado com %e foi por este motivo que esse %a relevância das contagens de palavras para classificação por perspectiva é patente% De acordo com a revisão, cinco doos seis artigos , cinco apresentaram um bom 

%Este capítulo apresentou duas hipóteses linguísticas assumidas por métodos baseados em frequências de palavras: 1) palavras específicas, denominadas \emph{banner words}, costumam ser utilizadas para defender perspectivas diferentes e 2)   %tão melhor quanto mais diferentes forem os empregos de palavras por perspectiva.

%ideia de que pessoas costumam utilizar palavras específicas, denominadas \emph{stigma words} ou \emph{banner words}, para defender perspectivas diferentes. O uso de classificadores baseados na presença ou frequência das palavras, portanto, costuma funcionar bem na identificação dessas perspectivas. Entretanto, em alguns casos, as construções do discurso refletem diferentes perspectivas melhor do que o uso de palavras específicas. Nestes casos, as taxas de acerto desses classificadores são menores do que o desejado. 

%Para ilustrar a relação entre as palavras de dois \emph{datasets} e o desempenho desses métodos, experimentos com o modelo de tópicos L-LDA foram executados. A extração das dez palavras mais fortemente associadas a cada tópico conduziu à visualização parcial de como o vocabulário dos \emph{datasets} se agrupa em torno de suas diferentes perspectivas. A informação, apesar de subjetiva, auxilia na compreensão das taxas de acerto obtidas com um classificador Naïve Bayes padrão, aplicado aos dois corpora. Para o primeiro \emph{dataset}, as taxas de acerto obtidas foram mais altas, o que pode ser explicado por uma presença maior de \emph{banner words} em comparação com o segundo \emph{dataset}.



% voa relação entre as palavras contidas em dois \emph{datasets} e suas diferentes perspectivas. No primeiro \emph{dataset}, as palavras extraídas para cada perspectiva se relacionam, semanticamente, melhor com a perspectiva em si do que no segundo \emph{dataset}. Neste último, palavras genéricas, como artigos e conjunções, se associam muito fortemente a todas as perspectivas, colaborando para uma uniformização do vocabulário. Isto explica, ainda que não seja trivial avaliar o quanto, as taxas de acerto mais altas obtidas com um Naïve Bayes no primeiro \emph{dataset}.

%Por fim, este capítulo também revisa artigos que discutem a questão do vocabulário do corpus. Estes artigos propõem metodologias para melhorar a classificação em \emph{datasets} cujo vocabulário é considerado muito uniforme - todas baseadas em algum esquema de citação a documentos. A questão do vocabulário do corpus é pouco discutida na literatura; provavelmente porque o problema é pouco comum e porque outros métodos, mais genéricos para \emph{datasets} difíceis de classificar, têm sido mais estudados.
 
%o primeiro \emph{dataset}, os artigos estão escritos sob uma das seguintes perspectivas: pró-Palestina ou pró-Israel. Por conta disso, o tópico \emph{pal} foi atribuído a todos os documentos pró-Palestina e o tópico \emph{isr}, a todos os pró-Israel. Estes tópicos facilitam a visualização das palavras mais fortemente associadas a cada uma das duas perspectivas, de acordo com o vocabulário empregado nos artigos. Um terceiro tópico, \emph{gen}, foi atribuído a todos os documentos, a fim de capturar as palavras que co-ocorrem neles independentemente de suas perspectivas. Após a execução do modelo, as 100 palavras mais fortemente associadas a cada um dos 3 tópicos foram coletadas. 35\% das associadas a \emph{isr} não estão presentes no conjunto recolhido para \emph{gen}. Para \emph{pal}, essa percentagem cai para 29\%. Por fim, 32\% das palavras mais fortemente associadas a \emph{isr} não fazem parte das 100 palavras mais fortemente associadas a \emph{pal} e vice-versa. Essas percentagens indicam que os autores dos documentos pró-Israel utilizam um vocabulário ligeiramente mais específico, na defesa de seus pontos de vista, do que os autores pró-Palestina. É importante ressaltar que nenhuma palavra foi filtrada na análise - ou seja, termos muito frequentes como \emph{the}, \emph{of} e \emph{and}, comumente extraídos dos \emph{datasets} antes da etapa de classificação, estavam presentes nos documentos processados pelo L-LDA.    

%Palavras associadas a \emph{isr} que não foram associadas a \emph{gen}
%['arafat', 'be', 'some', 'roadmap', 'us', 'yet', 'out', 'sharon', 'ariel', 'support', 'three', 'bush', 'palestine', 'new', 'terrorism', 'leader', 'then', 'jewish', 'after', 'arab', 'leadership', 'plan', 'president', 'than', 'bank', 'prime', 'regarding', 'like', 'could', 'violence', 'against', 'while', 'time', 'american', 'first']

%Palavras associadas a \emph{pal} que não foram associadas a \emph{gen}
%['then', 'some', 'authority', 'against', 'occupied', 'negotiations', 'occupation', 'sharon', 'united', 'end', 'way', 'palestine', 'international', 'be', 'after', 'plan', 'president', 'law', 'those', 'prime', 'land', 'i', 'violence', 'us', 'q', 'while', 'time', 'situation', 'first']

%Palavras associadas a \emph{pal} que não foram associadas a \emph{isr}
%['because', 'people', 'authority', 'states', 'right', 'occupied', 'any', 'negotiations', 'occupation', 'what', 'united', 'end', 'also', 'been', 'their', 'other', 'way', 'international', 'law', 'do', 'which', 'government', 'very', 'they', 'now', 'those', 'about', 'land', 'these', 'q', 'i', 'situation']

%Palavras associadas a \emph{isr} que não foram associadas a \emph{pal}
%['arafat', 'into', 'settlements', 'years', 'yet', 'out', 'even', 'would', 'ariel', 'west', 'support', 'three', 'bush', 'gaza', 'new', 'terrorism', 'leader', 'we', 'jewish', 'arab', 'most', 'leadership', 'minister', 'roadmap', 'than', 'bank', 'both', 'regarding', 'like', 'could', 'war', 'american']
%\textbf{As palavras estão ordenadas de acordo com a força da associação com cada um dos tópicos}


%\textbf{Lista de palavras e desempenho. Imagens e Tabelas.}

%Dado que em boa parte dos artigos estudados neste projeto, como \cite{lin-et-al2006} e \cite{klebanov}, atinge-se taxas de acerto superiores a 80\% com classificadores baseados em frequência de palavras, conclui-se que a mineração de perspectivas em discussões, artigos opinativos e debates requer metodologias diferentes, a depender de como as palavras foram escolhidas pelos autores dos documentos. %Nos debates estudados por \cite{hirst-et-al}, expressões de ataque e defesa são mais frequentes do que \emph{stigma words} e, como o método empregado no artigo foi um SVM treinado com frequências de palavras, observou-se que a classificação obtida para os lados do debate não refletia as perspectivas \emph{liberal} ou \emph{conservadora} - mas sim os lados \emph{oposição} (expressões de ataque) e \emph{situação} (expressões de defesa). 
%Estes estudos indicam a possibilidade de que, em debates e discussões nos quais há uma homogeneização do vocabulário empregado - o que pode acontecer quando todos os lados utilizam, em proporções similares, tanto expressões de ataque quanto de defesa -, classificadores baseados exclusivamente nas palavras utilizadas e/ou em suas frequências apresentarão má performance.

%A avaliação do desempenho desses classificadores\footnote{\textbf{SVMs e Naive Bayes padrão; LSPM}} nos \emph{datasets} estudados revela que artigos opinativos e notícias consolidaram perspectivas, através da escolha do vocabulário utilizado, melhor do que debates. Apesar disso, uma generalização neste sentido, restringindo o uso desses classificadores a artigos e notícias, não é recomendada por falta de indicativos linguísticos que comprovem essa tendência. Uma estratégia que pode ajudar na escolha ou descarte de um classificador desse tipo é uma análise das palavras que estão contidas nos documentos.

%\textbf{mostrar as percentagens prum dataset onde a linguagem eh mais uniforme. discutir.} 
%Explicar q isso nunca foi feito e q pessoas q ncontraram merda com essa feature fizeram outras coisas. descrever essas coisas.

%Estas hipóteses ilustram  
%Como indica \cite{}: 

%\textbf{MODO RASCUNHO AINDA}

%Explicar que uma diferença fundamental notada entre os datasets observados diz respeito à formalidade da linguagem empregada nos documentos. Enquanto alguns datasets utilizam documentos provenientes de meios onde a norma culta da linguagem impera, e uma revisão ortográfica é utilizada, outros são carregados de gírias, expressões e abreviações típicas da linguagem da Internet e contêm, eventualmente, grafias erradas para uma mesma palavra. <Dar exemplos textuais. Mostrar uma tabelinha, algo assim, indicando o volume de datasets com linguagem informal nos documentos estudados.> 

%A linguagem informal pode criar alguns desafios para a Mineração de Perspectiva [FECHAR O PROBLEMA EM: IDENTIFICAÇÃO DA PERSPECTIVA PRESENTE EM UM DOCUMENTO], especificamente no que diz respeito ao pré-processamento dos documentos e à escolha do método empregado. No pré-processamento, como indicado em (aaai-politics.pdf e 10.1.1.138.7160.pdf), a correção gramatical das palavras é bastante indicada. Com uma única versão de grafia para cada palavra (a correta), diminui-se a quantidade de ruído que grafias erradas podem causar na classificação.

%Uma característica dos datasets que deveria ser considerada sempre antes de se escolher o método utilizado - mas não é prática entre as pessoas que estudam perspective mining - é a frequência das palavras nos documentos do dataset. Se o léxico empregado pelos autores dos textos muda sensivelmente a depender de sua ideologia/perspectiva defendida/ponto de vista, é possível resolver o problema utilizando classificadores que usam essa frequencia como feature. GASTE TEMPO DANDO ALGUNS EXEMPLOS.  Em datasets informais, entretanto, como aponta Efron, a linguagem empregada por todos os lados da discussão pode ser basicamente a mesma: termos com forte carga de polaridade e gírias/jargões comuns na discussão [EXEMPLO]. Neste caso, é preciso utilizar métodos que utilizem mecanismos mais rebuscados do que a simples frequência/presença das palavras no texto. Alguns autores utilizam métodos mais gramaticais para lidar com datasets informais, como é o caso de , BLE e BLI. Os 2 primeiros criam o conceito de Opinion Frame (DEFINIR O CONCEITO). No primeiro caso, a ideia é associar corretamente opiniões a alvos, e associar alvos iguais utilizando técnicas de co-referência. Segundo Wiebe et al. (pegar a citação corretamente), ter as targets associadas, com as opiniões próximas, é uma boa forma de entender o overall de opiniões de um documento. A estratégia é uma alternativa possível para documentos em que as perspectivas estão muito associadas à linguagem opinativa, e onde essa linguagem é comum para todos os lados. Infelizmente, os resultados alcançados indicam que a tarefa não é trivial: blablablablababla (resolver coreferencia pra linkar targets não é tão simples assim => e dê exemplo). 

%uma diferença fundamental notada entre os \emph{datasets} observados diz respeito à formalidade da linguagem empregada nos documentos. Enquanto alguns datasets utilizam documentos provenientes de meios onde a norma culta da linguagem impera, e uma revisão ortográfica é utilizada, outros são carregados de gírias, expressões e abreviações típicas da linguagem da Internet e contêm, eventualmente, grafias erradas para uma mesma palavra. <Dar exemplos textuais. Mostrar uma tabelinha, algo assim, indicando o volume de datasets com linguagem informal nos documentos estudados.> 

%\begin{table}[t]
%\centering
%\begin{tabular}{| p{10cm} | }
%\hline

%\emph{"The recent \textbf{Israeli} government decision to begin building extensive walls
%around \textbf{Palestinian} is just one more example of how \textbf{Israeli} Prime
%Minister Ariel Sharon is unable to deal with \textbf{Israeli} problems save
%through his narrow security vision."} - Trecho extraído de artigo Pró-Palestina. \\ \hline

%\emph{"The first conclusion that the Israeli political and security
%establishment should learn and internalize after 18 months of
%\textbf{Palestinian} Intifada, concerns the intensity of \textbf{Palestinian} blind
%terrorism and guerilla warfare against the State of Israel."} - Trecho extraído de artigo Pró-Israel. \\ \hline

%\end{tabular}
%\label{3}
%\caption{Trechos com as palavras \emph{palestinian} e \emph{israeli}, extraídos do \emph{dataset} \textbf{Bitterlemons}.}
%\end{table}

%\begin{table}[t]
%\centering
%\begin{tabular}{| p{10cm} | }
%\hline

%\emph{"\textbf{Bush}
%and his advisers, who have been critical of Clinton's deep involvement
%in a failed peace process ever since taking office, nevertheless
%understood at the time that peace in the Middle East should be beyond
%politics in America, and that the US could not permit itself to turn its
%back on an Israeli leader who was determined to make peace."} - Trecho extraído de artigo Pró-Israel. \\ \hline

%\end{tabular}
%\label{4}
%\caption{Trecho com a palavra \emph{bush}, extraído do \emph{dataset} \textbf{Bitterlemons}.}
%\end{table}

%\begin{table}[t]
%\centering
%\begin{tabular}{| p{10cm} | }
%\hline
%\emph{"But just as we were close to a complete
%package that would have ended the \textbf{occupation} and established a
%Palestinian state, Barak permitted Ariel Sharon's provocative visit to
%Al Aqsa mosque, and launched his "revenge" on Palestinians."} - Trecho extraído de artigo Pró-Palestina. \\ \hline
%\end{tabular}
%\label{5}
%\caption{Trecho com a palavra \emph{occupation}, extraído do \emph{dataset} \textbf{Bitterlemons}.}
%\end{table}
% e  Termos associados ao Governo dos Estados Unidos, como \emph{bush} e \emph{american}, ilustram a relevância da aliança política estabelecida por Israel à época dos documentos. Palavras como \emph{occupation} e \emph{people}, por outro lado, ilustram a principal pauta Palestina no mesmo período: a ocupação da Faixa de Gaza.

%Os artigos deste primeiro \emph{dataset} foram escritos ou pelos editores do \emph{site} ou por convidados, divisão utilizada em \cite{lin-et-al2006} para avaliar o desempenho de um Naïve Bayes. No primeiro cenário, os documentos de treinamento eram os escritos pelos editores e os de teste, aqueles escritos pelos convidados; no segundo, tinha-se a situação inversa.

%As taxas de acerto obtidas com um classificador Naïve Bayes na identificação das perspectivas pró-Israel e pró-Palestina variaram entre 73.47\% e 98.98\%, a depender da divisão entre os conjuntos de treinamento e teste. %disponível em \cite{alibezz-nb}, as taxas de acerto obtidas foram de 73.47\% e 98.98\% para cada um dos cenários, respectivamente\footnote{A taxa de acerto obtida para o primeiro cenário foi significativamente inferior à obtida em \cite{lin-et-al2006} (85.85\%). Provavelmente, isto tem a ver com o número de iterações e as condições iniciais da execução.}. 

%O segundo \emph{dataset} estudado é o \textbf{Convote-Menor}, composto de colocações em debates da \emph{House of Representatives}, um dos dois órgãos principais do poder legislativo federal dos Estados Unidos. Os documentos foram marcados como sendo de parlamentares Republicanos ou Democratas, e como representando um posicionamento a favor ou contra a lei em pauta. Para este experimento, apenas a divisão entre Republicanos e Democratas foi considerada. O L-LDA foi aplicado a este \emph{dataset} de forma análoga ao primeiro experimento, e as dez palavras mais fortemente associadas a cada um dos tópicos - Genérico, Republicano e Democrata - estão listadas na Tabela 4.2. \emph{Stop words} também foram excluídas desta listagem.

%Ele é parte de um \emph{dataset} maior estudado pela primeira vez em \cite{thomas-pang-lee}, e está disponível em \cite{alibezz-convote}. A rotulação original dos documentos foi mantida neste \emph{dataset}: cada um deles é marcado com R, caso represente a colocação de um Republicano; ou D, no caso Democrata. Os documentos também são marcados com Y, caso representem a colocação de alguém que votou pela aprovação de uma lei; ou N, em caso contrário. Esta segunda marcação não foi explorada nesta monografia, nem a combinação das duas. O que se analisou, portanto, foi o desempenho do Naïve Bayes, disponível em \cite{alibezz-nb}, na classificação dos documentos como Republicanos ou Democratas.

%O L-LDA foi aplicado a este \emph{dataset} de forma análoga ao primeiro experimento, aproveitando os rótulos R e D, presentes nos documentos, como tópicos, e utilizando também um tópico genérico. As trinta palavras mais fortemente associadas a cada um dos tópicos, em ordem, podem ser conferidas na Tabela abaixo:


\chapter{Estudo de caso: Pontos de vista sobre o governo brasileiro}
\label{estudo}
%Onde se motiva o estudo das eleições e fala pq tratará de colunas de jornalistas de notoriedade nacional. TUDO NESTE CAPÍTULO PRECISA SER BEM DIDÁTICO. FALAR DA FOLHA DILMAFACTSBYFOLHA

%mtos trabalhos de mineração de perspectiva tratam de discussão política. aproveitando o fato de que em 2010 se realizam as primeiras eleições para presidente realmente acompanhadas pela Internet, a gente decidiu fazer um estudo de caso neste assunto. Há outras pessoas pensando em coisas parecidas, como a galera do eleitorando, q criou o serviço pra fazer identificação de opinião em tweets. mas oq estamos propondo é fazer um estudo de análise das perspectivas contidas em blogs e colunas de jornalistas de notoriedade nacional, por eles serem muito lidos e contribuírem para a formação de opinião de seus leitores sobre o governo brasileiro, oq tem tudo a ver com a hora da urna. Mostre depois a estrutura da parades. fale de outros corpus e eleitorando

%Por ser um ano eleitoral, a

Muitos dos trabalhos revisados neste projeto analisam documentos que tratam de política. Em particular, boa parte deles estuda textos relacionados a governos federais - quer sejam discussões entre os próprios governantes, como nos estudos de \textbf{Thomas, Pang e Lee} \cite{get-out-the-vote} ou \textbf{Hirst, Riabinin e Graham} \cite{hirst-et-al}, quer sejam artigos escritos por cidadãos ou especialistas, como nos artigos de \textbf{Mullen e Malouf}\footnote{Todos os trabalhos mencionados nesta sentença foram revisados no Capítulo \ref{chap3}.} \cite{aaai-politics} \cite{malouf-taking_sides}. Considerando essa tendência, e o fato de que 2010 é ano de eleições para presidente no Brasil, decidiu-se realizar um estudo de caso que aproveita a abundância de artigos, disponíveis na \emph{Web}, sobre o governo Lula e a sucessão presidencial. A ideia é construir um corpus com alguns desses documentos e investigar seus pontos de vista automaticamente, classificando-os de acordo com eles e analisando, de forma subjetiva, as palavras por eles enfocadas. É válido ressaltar que não foi encontrado nenhum outro trabalho, em português, que classifique documentos de acordo com seus pontos de vista.

%Os artigos escolhidos para este estudo foram extraídos de colunas, \emph{blogs} e \emph{sites} políticos mantidos por jornalistas de notoriedade nacional. A coleta de artigos contidos em \emph{blogs} escritos por cidadãos comuns também foi cogitada - entretanto, como eles são pouco conhecidos, comentados e divulgados, essa opção exigiria um esforço de análise manual dos \emph{posts} que foge ao escopo deste projeto. Além disso, uma vantagem em focar o estudo em material publicado por jornalistas conhecidos é poder correlacionar, posteriormente, os resultados obtidos a investigações sobre a formação de opinião na mídia brasileira \emph{online} - tanto na alternativa quanto na tradicional. A seleção dos veículos para este estudo de caso resultou do consenso entre a autora desta monografia e dois jornalistas \textbf{COMO CITÁ-LOS, DIZER SEUS NOMES?}. O critério básico para as escolhas foi a defesa clara de um ponto de vista sobre o governo Lula e/ou a sucessão presidencial de 2010. 

As próximas seções deste capítulo se estruturam da seguinte forma: na seção \ref{estudo:sec1}, a construção do corpus é apresentada - desde a seleção dos veículos até a definição dos pontos de vista explorados; na seção \ref{estudo:sec2}, experimentos com um classificador Naïve Bayes\footnote{Esse classificador é apresentado no Capítulo \ref{basicos}.} são conduzidos para, assim como em outros trabalhos revisados neste projeto, se classificar artigos de acordo com seus pontos de vista; na seção \ref{estudo:sec3}, o modelo de tópicos L-LDA\footnote{Esse modelo é apresentado no Capítulo \ref{basicos}.} é aplicado ao corpus, evidenciando o uso de palavras por artigos com posicionamentos diferentes. 

% posicionamentos é analisada  
%Por serem muito lidos, e defenderem perspectivas claras nos textos que publicam, estes jornalistas têm o poder de influenciar muitos eleitores brasileiros. %análises  % mas requereria muito mais esforço de análise do que o uso exclusivo de documentos  %Lula é minha anta Dois jornalistas

\section{Construindo um corpus para estudo}
\label{estudo:sec1}

%O corpus deste estudo é composto de artigos escritos para colunas de jornais e revistas ou para \emph{sites} sobre política. São textos de caráter eminentemente opinativo, que diferem de notícias tanto pela finalidade - \textbf{blibli} - quanto pela linguagem - \textbf{blabla}. 

Os artigos escolhidos para este estudo foram extraídos de colunas, \emph{blogs} e \emph{sites} políticos mantidos por jornalistas de notoriedade nacional. A coleta de \emph{posts} de \emph{blogs} escritos por cidadãos comuns também foi cogitada - entretanto, como eles são pouco conhecidos, comentados e divulgados, essa opção exigiria um esforço de análise manual dos \emph{posts} que foge ao escopo deste projeto. Além disso, uma vantagem em focar o estudo em material publicado por jornalistas conhecidos é poder correlacionar, posteriormente, os resultados obtidos a investigações sobre a formação de opinião na mídia brasileira \emph{online} - tanto na alternativa quanto na tradicional. 

A seleção dos veículos para este estudo de caso resultou do consenso entre a autora desta monografia e os jornalistas Lucas Cunha Almeida e Andrea Duarte Bessa\footnote{Lucas Cunha, atualmente, é repórter do jornal baiano A Tarde; Andrea Bessa é gerente de comunicações da empresa Cetrel.}. O critério básico para as escolhas foi a defesa clara de um ponto de vista sobre o governo Lula e/ou a sucessão presidencial de 2010. Assim como em outros artigos revisados nesta monografia, que dividem os corpora analisados em dois lados antagônicos, assume-se que os artigos do \emph{dataset} desse estudo de caso dividem-se entre pró-situação e pró-oposição. Neste sentido, há autores que apóiam candidatos de ambos os lados, muitas vezes criticando o governo Lula e suas personalidades. Os trechos de \emph{posts} apresentados nessa seção reforçam a ideia de que essa divisão é adequada. O lado pró-situação é composto de artigos veiculados em:


%EXPLICAR Q POR SIMPLICIDADE E PQ QUASE TUDO ESTUDADO É ASSIM, DIVIDIU-SE AS PERSPECTIVAS A SEREM IDENTIF. AUTOM. EM PRÓ E ANTI GOVERNO
\begin{enumerate}

\item \textbf{Luis Nassif Online}\footnote{http://www.advivo.com.br/luisnassif/} Este é o \emph{blog} do jornalista Luis Nassif, premiado como Melhor Blog de Política pelo iBest 2008\footnote{http://idgnow.uol.com.br/internet/2008/05/21/ibest-2008-anuncia-vencedores/}. Nassif, que já trabalhou na TV Cultura e Rede Bandeirantes, mantém o \emph{blog} há cinco anos, enfocando principalmente assuntos relativos à política brasileira. Artigos do \emph{blog} são frequentemente citados, de forma positiva, em veículos de campanha pró-situação, como os \emph{sites} Blog da Dilma\footnote{http://dilma13.blogspot.com/} e Os Amigos do Presidente Lula\footnote{http://osamigosdopresidentelula.blogspot.com/}. De fato, o Luis Nassif Online adota um posicionamento pró-situação, como comprovam os trechos a seguir:

\begin{quote}
\emph{"Desde o ano passado, estava claro [sic] a falta de competitividade de José Serra, seja por não ter feito um governo brilhante em São Paulo, por não representar o novo e por não conseguir desenvolver um discurso próprio."}

{\small Retirado de \emph{"Em Minas, a mãe de todas as batalhas"} - 02/09/2010} 
\end{quote}

\begin{quote}

\emph{"Na entrevista, Bonner se limitou a perguntar da dependência de Dilma em relação à Lula [...] A consequência foi Dilma rebatendo com facilidade cada bobagem dita, reforçando o discurso social, mas sem avançar em uma proposta sequer de programa, explicando a lógica das alianças políticas. E William Bonner interrompendo-a a toda hora, impedindo sequer uma resposta completa. Algo tão desastrado e mal educado que obrigou Fátima Bernardes, do alto de sua elegância, a calá-lo com um sinal, para que parasse de ser inconveniente."}

{ \small Retirado de \emph{"O dia em que William Bonner escorregou"} - 10/08/2010}

\end{quote}

%Além de conter artigos escritos pelo próprio Luis Nassif, o \emph{blog} também publica textos de outros autores - alinhados, evidentemente, com os pontos de vista do jornalista. 

Como o veículo possui muito conteúdo, foram considerados apenas os artigos da categoria "Eleições".  

\item \textbf{Conversa Afiada}\footnote{http://www.conversaafiada.com.br/} O \emph{site} se define como um portal de jornalismo independente, contendo principalmente artigos produzidos por Paulo Henrique Amorim. O jornalista, que já trabalhou para as Redes Globo e Bandeirantes e para a revista Carta Capital, mantém o \emph{site} desde 2006. Enfocando a política brasileira, o Conversa Afiada apóia, dentre outras iniciativas do governo federal, a candidatura da ex-ministra Dilma Rousseff\footnote{http://www.conversaafiada.com.br/brasil/2010/07/02/mino-explica-por-que-apoia-a-dilma-porque-ela-e-melhor-que-o-serra/}. Os trechos abaixo justificam a escolha do \emph{site} como representante da mídia \emph{online} pró-situação:

\begin{quote}

\emph{"O Governo Lula é um sucesso e a popularidade dele, recordista desde o primeiro dia de Governo. Promoveu a inclusão social, ampliou a classe média e assistiu os pobres. Fez uma política externa que não tirou o sapato para os Estados Unidos. A Dilma é a sua legítima sucessora: foi a CEO do Governo Lula. O Serra é um nada."}

{\small Retirado de \emph{"A Dilma não é um tsunami. Dilma é o rio que segue para o mar"} - 27/08/2010}
\end{quote}

\begin{quote}

\emph{"Segundo a tevê DEMO-Tucana da Bahia, a afiliada da Globo, Jacques Wagner está na frente de Paulo Souto por 46\% a 19\%. Paulo Souto é o aliado de Serra na Bahia. A TV Bahia, também."}

{\small Retirado de \emph{"Sumiram com o dinheiro do Serra. Serra é barrado em procissão"} - 07/08/2010}
\end{quote}

%Os artigos do \emph{site} muitas vezes citam notícias retiradas de portais, como o G1 ou UOL, a fim de contextualizar críticas e análises.

Também por possuir muito conteúdo, apenas os artigos pertencentes à categoria "Política" foram considerados.
% Prêmio iBest de Melhor Blog de Política, em eleição popular e da Academia iBest. 

\item \textbf{Escrevinhador}\footnote{http://www.rodrigovianna.com.br/} O \emph{blog}, mantido pelo \emph{site} da revista Caros Amigos, é escrito pelo jornalista Rodrigo Vianna, que também é repórter da Rede Record. Ele está no ar desde 2008, enfocando acontecimentos da vida política do Brasil e do Mundo. No que diz respeito ao Brasil, o conteúdo do \emph{blog} assume uma perspectiva pró-situação, como ilustram os trechos abaixo:

\begin{quote}

\emph{"Abandonado pelos aliados do DEM e do PSDB, em queda nas pesquisas, Serra refugia-se na mídia. O candidato do PSDB virou isso: porta-voz dos interesses da velha mídia. Faz sentido. É quem, em última instância, sustenta a candidatura."}

{\small Retirado de \emph{"Serra, porta-voz da velha mídia; é Zé ou Mané?"} - 19/08/2010}
\end{quote}

\begin{quote}
\emph{"O programa da Dilma foi um show. [...] Foi um programa em que Lula não apareceu mais que Dilma, e nem sumiu – porque seria falso, ela é a candidata dele. Foi um programa em que Lula passou o bastão a Dilma. De forma eficiente, corajosa e, ao mesmo tempo, emocionante."}

{\small Retirado de \emph{"Dilma acerta a mão; Serra quer virar 'Zé'"} - 18/08/2010}

\end{quote}

Como o \emph{blog} também trata de outros assuntos, apenas as categorias "Plenos Poderes" e "Palavra Minha", mais direcionadas à política, foram consideradas para extração de artigos.

\item \textbf{Brasília, eu vi}\footnote{http://brasiliaeuvi.wordpress.com/} O \emph{blog}, escrito pelo jornalista Leandro Fortes, que também trabalha para a revista Carta Capital, agrega alguns de seus artigos para a revista e outros textos sobre política. Estes artigos têm boa recepção em \emph{sites} de campanha pró-situação, como o Blog da Dilma\footnote{http://dilma13.blogspot.com/2010/08/caso-lunus-verdade-dos-fatos.html}. De fato, eles assumem uma perspectiva de defesa da situação, como justificam os trechos abaixo:


\begin{quote}
\emph{"Assim, enquanto a imprensa mundial se dedica a decodificar as engrenagens e circunstâncias que fizeram de Lula o mais importante líder mundial desse final de década, a imprensa brasileira se debate em como destituí-lo de toda glória, de reduzí-lo a um analfabeto funcional premiado pela sorte, a um manipulador de massas movido por programas de bolsas e incentivos [...]."}

{\small Retirado de \emph{"Não verás Lula nenhum"} - 18/05/2010}
\end{quote}

\begin{quote}

\emph{"Ao acusar o presidente Luiz Inácio Lula da Silva de ter transformado o Brasil em uma “república sindicalista”, José Serra optou por agregar a seu modelito eleitoral, definitivamente, o discurso udenista de origem, de forma literal, da maneira como foi concebido pelas elites brasileiras antes do golpe militar de 1964."}

{\small Retirado de \emph{"Serra precisa de mais amigos"} - 15/07/2010}
\end{quote}
\end{enumerate}


O lado pró-oposição, por sua vez, é composto de artigos veiculados em:

\begin{enumerate}

\item \textbf{Reinaldo Azevedo}\footnote{http://veja.abril.com.br/blog/reinaldo/} O \emph{blog}, escrito pelo jornalista homônimo, é mantido pela revista Veja. Autor da frase \emph{"Tudo que é bom para o PT é ruim para o Brasil"} \cite{bom-pt-mau-brasil}, Reinaldo Azevedo, que já foi editor da Folha de S. Paulo, alimenta seu \emph{blog} com críticas ao governo atual, como evidenciam os trechos abaixo:

\begin{quote}

\emph{"O problema dos petistas é que eles são viciados no aulicismo, na cortesania. Ao conviver com pessoas que sempre têm um preço, ficam chocados e tomam como ofensa pessoal a descoberta de que nem todos se comportam com essa moral anã."}

{\small Retirado de \emph{"Presidente do PT repete ladainha autoritária do programa 'Rubriquei, mas não traguei'. Ou: 'Ai que vontade de censurar a Veja!!!' Contenha a coceira, companheiro!} - 15/07/2010} 
\end{quote}

\begin{quote}

\emph{"Cinco centrais sindicais assinaram um vergonhoso manifesto contra a candidatura do tucano José Serra à Presidência. Antes de mais nada, e a despeito da mentira essencial que está contida no texto — já falo a respeito —, cumpre destacar: trata-se de um manifesto ilegal, de mais um crime eleitoral escancarado."}

{\small Retirado de \emph{"Acuado pelo 'Rubriquei, mas não traguei', PT mobiliza centrais sindicais. E elas assinam um documento ilegal e mentiroso."} - 12/07/2010 }

\end{quote}

\item \textbf{Coluna do Augusto Nunes}\footnote{http://veja.abril.com.br/blog/augusto-nunes/} A coluna, parte da revista Veja, é escrita pelo jornalista Augusto Nunes, que também apresenta o programa Roda Viva na TV Cultura. Seus artigos têm má recepção em alguns veículos que defendem o atual governo, como o Luis Nassif Online\footnote{http://www.advivo.com.br/blog/luisnassif/serra-e-fhc-uma-relacao-delicada} e o Blog da Dilma\footnote{http://dilma13.blogspot.com/2010/01/mais-uma-do-tucano-augusto-nunes.html}, justamente por assumirem uma posição pró-oposição. Os trechos abaixo justificam esta perspectiva:


\begin{quote}

\emph{"Como todo sinal de alarme, o som de um neurônio em ebulição é  perturbador, mas muito útil. Quem tem juízo entenderá que Dilma Rousseff não é uma candidata em campanha. É uma ameaça a caminho."}

{\small Retirado de \emph{"O som perturbador do neurônio em ebulição"} - 20/07/2010}

\end{quote}

\begin{quote}

\emph{"O eleitor merece saber se Lula recebeu uma herança maldita e reconstruiu o país, como repete há pelo menos seis anos, ou se resolveu valer-se de mentiras e fantasias para desqualificar o legado do antecessor que acabou com a inflação, consolidou a democracia constitucional e fixou diretrizes econômicas que, em sua essência, vigoram até hoje."}

{\small Retirado de \emph{"FHC aceita o convite para o duelo que Lula não pode recusar."} - 11/02/2010}
\end{quote}

Todos os artigos extraídos dessa coluna pertencem à categoria "Direto ao Ponto", por ela tratar especificamente da política brasileira atual.

\item \textbf{Coluna do Diogo Mainardi}\footnote{http://veja.abril.com.br/blog/mainardi/} A coluna, escrita desde 2002, é a mais lida da revista Veja segundo ela mesma, reunindo críticas à política e à economia brasileiras. O jornalista Diogo Mainardi se opõe aos governos petistas, tendo inclusive publicado, em 2007, o livro Lula é Minha Anta \cite{lula-anta}, no qual agrupa diversos artigos escritos para sua coluna na Veja. Os trechos abaixo ilustram a posição de Mainardi como um grande crítico do governo do PT e de sua candidata Dilma Rousseff:

\begin{quote}

\emph{"Dilma Rousseff teve uma loja de produtos importados. O empreendimento durou menos de um ano e meio. Se Dilma Rousseff mostrar como presidente da República o mesmo talento que mostrou como empresária, o Brasil já pode ir fechando as portas."}

{\small Retirado de \emph{"Dilma 1,99 Rousseff"} - 04/09/2010}
\end{quote}  

\begin{quote}

\emph{"No futuro, quando alguém quiser relatar os fatos deste período, terá de recorrer necessariamente aos processos judiciais, que detalharam o modo lulista de se organizar, de se acumpliciar, de se infiltrar e de fazer negócios."}

{\small Retirado de \emph{"A história em inquéritos"} - 20/03/2010}
\end{quote}

\item \textbf{Portal de Carlos Alberto Sardenberg}\footnote{http://www.sardenberg.com.br/site/index.php} O portal contém artigos do jornalista para suas colunas nos jornais O Globo e O Estado de S. Paulo, além de outros textos de análise política e econômica. Além destas ocupações, Sardenberg também é comentarista da TV Globo e âncora da Rádio CBN, tecendo comentários sobre a economia mundial e brasileira. Os trechos abaixo transparecem seu posicionamento pró-oposição: 

\begin{quote}

\emph{"O governo Lula não quer fazer concessões à iniciativa privada porque está num ímpeto estatizante, em ano eleitoral. Só que o Estado não tem os recursos para fazer nada de substancial. Fica por isso mesmo."}

{\small Retirado de \emph{"As tarefas de Lula"} - 22/03/2010}
\end{quote}

\begin{quote}

\emph{"É verdade que o país está de novo em um bom momento. Mas não é verdadeira a conclusão que o 'lulismo' tira disso: que isso tudo só está acontecendo porque Lula é o presidente."}

{\small Retirado de \emph{"A salvação?"} - 01/04/2010}
\end{quote}
\end{enumerate}

É válido ressaltar que os autores dos artigos muitas vezes colocam trechos de notícias, ou mesmo textos críticos de outros autores, em seus escritos, colaborando para a riqueza da linguagem no corpus.

Outros veículos foram cogitados, como o Blog do Noblat\footnote{http://oglobo.globo.com/pais/noblat/}, o \emph{blog} de Miriam Leitão para o jornal O Globo\footnote{oglobo.globo.com/economia/miriam/}, a coluna de Cristiana Lôbo para o portal G1\footnote{http://g1.globo.com/platb/cristianalobo/} e o \emph{blog} de Celso Ming para o jornal O Estado de S. Paulo\footnote{blogs.estadao.com.br/celso-ming/ }. Os posicionamentos contidos nestes veículos, entretanto, não foram considerados claros o suficiente para os propósitos deste estudo. 

Todos os artigos contidos nas colunas, \emph{sites} e \emph{blogs} selecionados foram  publicados entre 01/01/2010 e 06/09/2010. O período fixado, por fazer parte de um ano eleitoral, encerra uma quantidade significativa de artigos pró-situação e pró-oposição. Por este motivo, e também para manter o escopo do estudo atrelado às eleições 2010, artigos de anos anteriores não foram coletados. A extração dos documentos foi feita de forma automatizada com \emph{scripts} desenvolvidos nas linguagens de programação Python\footnote{http://python.org/} e UNIX ShellScript\footnote{http://www.gnu.org/software/bash/manual/bashref.html}. Todos eles estão disponíveis no repositório \emph{online} de Aline Bessa\footnote{http://github.com/alibezz}. Como os jornalistas eventualmente publicam sobre política mundial ou outros assuntos, foi feita uma filtragem nos artigos, de modo a restarem apenas aqueles que contêm pelo menos uma das seguintes palavras-chave: "Lula", "FHC", "Dilma", "Serra", "Marina", "PT", "PV", "PSDB". Todos os documentos foram, por fim, anonimizados, para que os nomes de seus autores não interferissem nos estudos. 

Após filtragem e anonimização, restaram 1065 artigos pró-situação e 2747 pró-oposição. Para os estudos feitos com o corpus, envolvendo o classificador Naïve Bayes e o modelo de tópicos L-LDA, reduziu-se a quantidade de documentos pró-oposição para 1065, utilizando-se apenas 550 dos 2377 artigos extraídos do \emph{blog} de Reinaldo Azevedo, amostrados aleatoriamente. Essa estratégia foi adotada porque o classificador se mostrou sensível a quantidades muito discrepantes de palavras por ponto de vista. Como o uso do L-LDA estende as análises feitas com o classificador, decidiu-se manter o corpus idêntico para ambos os estudos. 

\begin{table}[t]
\centering
\begin{tabular}{| l | l | p{5cm} | }
\hline

\textbf{Veículo} & \textbf{Coleta} & \textbf{Filtragem/Anonimização} \\ \hline

\textbf{Reinaldo Azevedo} & 2490 & 2377\ensuremath{^*} \\ \hline
\textbf{Augusto Nunes} & 579 & 450 \\ \hline
\textbf{Diogo Mainardi} & 40 & 32 \\ \hline
\textbf{Carlos Sardenberg} & 59 & 33  \\ \hline
\textbf{Conversa Afiada} & 375 & 337  \\ \hline
\textbf{Luis Nassif Online} & 994 & 525 \\ \hline
\textbf{Escrevinhador} & 222 & 179  \\ \hline
\textbf{Brasília, eu vi} & 34 & 24  \\ \hline
\end{tabular}
\label{tab1:estudo}
\caption{Quantidades de artigos disponíveis em cada etapa da construção do corpus. \ensuremath{^*}Apenas 550, amostrados aleatoriamente, foram aproveitados.}
\end{table}

O número de artigos coletados em cada veículo varia bastante, como pode ser observado na Tabela 4.1. No corpus sobre o conflito Israel-Palestina, estudado por Lin et al. em trabalho revisado no Capítulo \ref{chap3}, este comportamento também foi observado. Os documentos desse \emph{dataset} foram escritos por diversos autores, e alguns deles contribuíram com muito mais textos do que outros. Ainda assim, como pode ser verificado na revisão apresentada no Capítulo \ref{chap3}, os resultados obtidos são de alta qualidade \cite{lin-et-al2006}. Isto reforça a ideia de que essa variação não interfere significativamente na qualidade dos experimentos feitos com o corpus deste estudo de caso.


\section{Classificando documentos com um Naïve Bayes}
\label{estudo:sec2}

O primeiro estudo conduzido com esse corpus consiste na classificação dos artigos de acordo com os pontos de vista pró-oposição e pró-situação. Esse capítulo não se propõe a comparar classificadores diferentes, de modo que o Naïve Bayes será a única técnica explorada em seu escopo. O classificador Naïve Bayes foi escolhido porque, de acordo com a revisão do Capítulo \ref{chap3}, não há um consenso sobre qual dos dois - ele ou um SVM - é mais indicado para classificação por ponto de vista. Em alguns trabalhos o Naïve Bayes conduz a melhores resultados; em outros, o SVM. Essas questões, e a consideração de que o Naïve Bayes é mais simples de implementar, justificam sua escolha.  

O classificador explora apenas as contagens de palavras dos documentos para definir suas classes. Os motivos para essa escolha foram: popularidade do uso nos trabalhos revisados no Capítulo \ref{chap3}, o que indica essa escolha como a mais natural; simplicidade de computação, dado que essa escolha envolve apenas a soma das ocorrências de palavras; fácil adaptação a qualquer língua, em particular a portuguesa. 

% por contagens de palavras também são os mesmos apontados no Capítulo \ref{chap3}, na seção \ref{freqs:experim}. 

Essas escolhas se mostraram adequadas para o problema: a taxa de acerto obtida na classificação foi de 89.43\%. Assim como em outros artigos revisados no Capítulo \ref{chap3} \cite{lin-et-al2006} \cite{aaai-politics} \cite{klebanov}, esse valor foi obtido via validação cruzada de dez dobras\footnote{Essa técnica de validação é apresentada no Capítulo \ref{basicos}.}. A precisão, rechamada e métrica F1\footnote{Essas métricas são descritas no Capítulo \ref{basicos}.} associadas a cada classe (ponto de vista) constam na Tabela 4.2. %tabela mais precisa %precisao 0 0,915979712 precisao 1 0,877579312 % rechamada 0 0,870754717 1 0,917924528  %f1 0 0,891837888 1 0,896484354

%; a precisão média, de  89.68\%; a métrica F1 média, de 89.42\%

\begin{table}[h]
\centering
\begin{tabular}{| l | l | l | l | }
\hline

\textbf{Classe} & \textbf{Precisão} & \textbf{Rechamada} & \textbf{Métrica F1} \\ \hline

\textbf{Pró-situação} & 91.60\% & 87.08\% & 89.18\% \\ \hline
\textbf{Pró-oposição} & 87.78\% & 91.79\% & 89.65\% \\ \hline
\end{tabular}
\label{resultados}
\caption{Métricas de desempenho associadas a cada classe.}
\end{table}

A primeira coluna da Tabela 4.2 indica que o Naïve Bayes classificou menos documentos erroneamente como pró-situação do que como pró-oposição. Em outras palavras, a classificação como pró-situação foi mais \textbf{precisa}. A segunda coluna indica que, por outro lado, o Naïve Bayes classificou mais documentos pró-oposição corretamente. Neste sentido, a rechamada funciona como a taxa de acerto. Por fim, a terceira coluna indica que, ponderando as colunas anteriores via métrica F1, os resultados obtidos são praticamente idênticos. Isso significa, na prática, que o Naïve Bayes apresenta um desempenho bastante equilibrado considerando-se ambas as classes. 

O bom desempenho obtido com o Naïve Bayes indica que o uso de metodologias mais complexas, envolvendo aspectos sintáticos dos documentos, por exemplo, não é necessário. Além disso, é válido ressaltar que a adoção de características sintáticas/semânticas, para a classificação de um corpus em português, não é tão simples como para um corpus em inglês. Isso advém  da carência de \emph{softwares} e dicionários semânticos para a língua portuguesa - em particular, a variação utilizada no Brasil. Esses \emph{softwares} e dicionários colaboram com a determinação de diversos aspectos gramaticais/linguísticos dos textos, minimizando o esforço manual requerido quando eles não podem ser utilizados.


Os valores obtidos com o classificador Naïve Bayes são comparáveis àqueles obtidos por \textbf{Durant e Smith} em seu trabalho sobre George W. Bush e a Guerra do Iraque\footnote{Esse trabalho foi revisado no Capítulo \ref{chap3}.} \cite{durant-smith}: 89.77\%. É válido ressaltar que, diferentemente da metodologia adotada por Durant e Smith, nenhuma palavra foi descartada no processamento dos textos para este estudo de caso.  

Os artigos escolhidos para este corpus são compostos, muitas vezes, de textos de outros autores. Isto reforça a ideia de que o classificador Naïve Bayes está efetivamente aprendendo os pontos de vista dos documentos, em vez de estilos de escrita. De todo modo, assim como no estudo de \textbf{Lin et al.} sobre o conflito Israel-Palestina\footnote{Esse trabalho foi revisado no Capítulo \ref{chap3}.} \cite{lin-et-al2006}, foi conduzido um experimento em que os artigos pertencentes aos conjuntos de treinamento e teste são escritos por autores diferentes. Se o que está sendo aprendido são de fato os pontos de vista dos documentos, o desempenho do classificador não deve ser muito diferente do obtido na validação cruzada de dez dobras. Testando com artigos da coluna de Augusto Nunes e do \emph{site} Conversa Afiada, e treinando com os demais, a taxa de acerto obtida foi de 92.79\%, acompanhada de precisão média de 93.32\% e métrica F1 média de 91.86\%. Este experimento, portanto, ratifica os outros resultados, evidenciando que o classificador Naïve Bayes cumpre bem a tarefa de identificar os pontos de vista pró-oposição e pró-situação.

% indicado para identificar perspectivas em artigos opinativos.
%Uma alta taxa de acerto também indica que a identificação de perspectiva política em \emph{blogs} brasileiros, cujas linguagens se assemelham às estudadas, pode apresentar um bom resultado.

%Se a taxa de acerto for alta, é possível pensar em aplicações que envolvam a identificação da perspectiva d


%- talvez até escritas por blogueiros comuns, em vez de jornalistas.


\section{Ilustrando o uso de palavras por ponto de vista}
\label{estudo:sec3}
%Onde se explica o uso do L-LDA, o overlap de palavras, mostra-se snippets e talz.

A análise das principais palavras enfocadas por cada ponto de vista amplia a compreensão dos resultados obtidos na seção anterior. A alta taxa de acerto obtida sugere que os jornalistas se expressam de forma bastante diferente, de modo que o Naïve Bayes não \emph{se engana} com muita frequência. Em outras palavras, como essas expressões foram quantificadas via contagens de palavras, a alta taxa de acerto indica que essas características foram suficientes para discrimar documentos de classes diferentes. Essa seção se propõe a ilustrar quais palavras foram mais destacadas por eles, colaborando para a transmissão de pontos de vista diferentes.

Para a investigação sobre essas palavras, foi feita uma aplicação do modelo L-LDA. Cada documento foi associado a dois tópicos: um neutro, igual para todos eles, e um referente a seu ponto de vista (pró-oposição ou pró-situação). Há, portanto, três tópicos diferentes nesta aplicação. O uso de um tópico neutro, associado a todos os documentos, ajuda a filtrar palavras muito comuns nos \emph{datasets}, independentemente de ponto de vista. Dois exemplos são as palavras \emph{vez} e \emph{campanha}. Essa é a diferença fundamental entre essa aplicação do L-LDA e a simples contagem de palavras em documentos, dividida entre os dois pontos de vista. Esse tipo de contagem não diferencia quais palavras são \emph{mais destacadas} em documentos escritos sob uma certa perspectiva e quais são \emph{muito utilizadas} por todos eles. Essas últimas possivelmente não os discriminam tanto quanto as primeiras.

% Cada documento de ambos os \emph{datasets} foi associado a dois tópicos: um neutro, idêntico para todos eles, e outro referente a seu ponto de vista. No primeiro corpus, esses pontos de vista são pró-Israel ou pró-Palestina; no segundo, republicano ou democrata. Há portanto, em cada corpus, três tópicos diferentes.


%O motivo para isso é  que a classificação em si se baseia nesses usos de palavras, através de suas contagens. Neste sentido, portanto, a compreensão de quais palavras receberam mais destaque 

% o modelo generativo L-LDA foi aplicado ao corpus. Cada artigo O uso desses tópicos evidencia quais palavras são mais enfocadas por cada ponto de vista. % e quais são utilizadas de forma muito semelhante em todo o corpus (genérico).

%Para a investigação sobre o uso de palavras,
A Tabela 4.3 elenca as dez palavras mais frequentemente associadas a cada ponto de vista. O tópico neutro não será analisado, haja vista que foi utilizado apenas para auxiliar na filtragem das palavras mais comuns do corpus. Em alguns casos, essas palavras coincidem nos dois pontos de vista. Isso significa que ambos deram muito destaque a elas.

%em dois ou mais tópicos. Se o genérico não estiver presente, i

%essa palavra. Em caso contrário, tem-se que a palavra, além de ser muito enfocada por pelo menos uma das duas perspectivas, é também uma das mais comuns no corpus.

\begin{table}[h]
\centering
\begin{tabular}{| l | p{10cm} | }
\hline
\textbf{Tópico} & \textbf{Palavras} \\ \hline
\textbf{Pró-situação} & serra, dilma, lula, psdb, presidente, pt, candidato, folha, tucano, eleições\\ \hline
\textbf{Pró-oposição} & dilma, lula, presidente, brasil, rousseff, pt, gente, candidata, mundo, candidato \\ \hline
\end{tabular}
\label{tab:palavras}
\caption{As dez palavras mais frequentemente associadas aos tópicos pró-situação e pró-oposição, em ordem e excluindo-se artigos, preposições, conjunções, advérbios e pronomes pessoais.}
\end{table}

%\textbf{Genérico} &  governo, brasil, serra, estado, lula, poder, presidente, nacional, vez, campanha \\ \hline

%JUSTIFICAR ISSO DIREITO, E FAZER AS FIGURAS
A Tabela 4.3 indica que as palavras enfocadas por pontos de vista diferentes são as mesmas em alguns casos, diferindo apenas na forma como são enfatizadas. Em outras palavras, a ordem das palavras nessa tabela corresponde, diretamente, ao quanto elas se associaram a cada tópico. Os artigos pró-oposição, por exemplo, dão muito destaque às palavras \emph{lula} e \emph{dilma}; os pró-situação, por sua vez, também enfatizam estas palavras, mas dão um destaque maior a \emph{serra}, candidato à presidência pelo PSDB. A associação de palavras semelhantes aos tópicos pró-oposição e pró-situação advém do fato de que os artigos compartilham um tema geral - o governo brasileiro - e, consequentemente, o mesmo vocabulário básico. É diferente do que acontece quando os tópicos correspondem a temas diferentes em vez de pontos de vista, como pode ser visto no trabalho de Ramage et al. sobre \emph{tags} de \emph{blogs} \cite{llda}.

Para visualizar melhor a ênfase dada às palavras na Tabela 4.3, elas foram processadas pelo \emph{software} wordle\footnote{http://wordle.net}, resultando nas figuras \ref{situacao} e \ref{oposicao}. O tamanho das palavras nas imagens está relacionado ao número de vezes que elas se associam a cada tópico. As imagens evidenciam o destaque dado aos políticos Lula, Dilma Rousseff e José Serra nos artigos analisados. %A imagem \ref{generico} dá certo destaque a Lula e José Serra, mas também enfatiza outros termos, como \emph{governo} e \emph{brasil}, relacionados mais genericamente à temática dos artigos: a política brasileira.

%Isso fica mais claro nas Figuras \ref{situacao} e \ref{oposicao}. É válido ressaltar que, nos experimentos do Capítulo \ref{chap3}, foram geradas imagens como essas, mas elas não indicaram ênfases muito diferentes no uso das palavras selecionadas.
\begin{figure}[h]
  \centering % este comando é usado para centralizar a Figura
  \includegraphics[width=12.5cm, height=5cm]{situacao.png}\\
  \caption{Representação gráfica para as dez palavras associadas ao tópico pró-situação na Tabela 4.3.}
  \label{situacao}
\end{figure}

\begin{figure}[h]
  \centering % este comando é usado para centralizar a Figura
  \includegraphics[width=11.5cm, height=6.5cm]{oposicao.png}\\
  \caption{Representação gráfica para as dez palavras associadas ao tópico pró-oposição na Tabela 4.3.}
  \label{oposicao}
\end{figure}



%\begin{figure}[h]
%  \centering % este comando é usado para centralizar a Figura
%  \includegraphics[width=12.5cm, height=6.5cm]{generico.png}\\
%  \caption{Representação gráfica para as dez palavras associadas ao tópico genérico na Tabela 4.3.}
%  \label{generico}
%\end{figure}

Essas figuras também sugerem que os artigos pró-situação dão mais enfoque a personalidades relacionadas à situação, como Lula e Dilma Rousseff, do que os pró-oposição a personalidades da oposição, como José Serra ou Marina Silva. Essa última, inclusive, candidata à presidência pelo PV, não foi mencionada nas palavras listadas na Tabela 4.3. Essas duas observações evidenciam que os veículos, no período analisado, concentraram seus antagonismos em personalidades políticas dos partidos PT, como Lula e Dilma Rousseff, e PSDB, como José Serra. 

%É válido ressaltar, por fim, que, apesar dos textos terem caráter crítico e não necessariamente informativo, as palavras elencadas na Tabela 4.3 não carregam uma polaridade natural, como no caso dos adjetivos \emph{bom} ou \emph{ruim}. Isso reforça a ideia de que a expressão de pontos de vista não é tão simples quanto a emissão de opiniões pontuais e claramente polarizadas, como \emph{Serra é um candidato ruim} ou \emph{Lula é um presidente bom}. 

Os pontos de vista desse corpus são construídos através de argumentações complexas, como indicam os trechos de artigos apresentados na seção \ref{estudo:sec1}. Para compreender melhor a relação que as palavras da Tabela 4.3 estabelecem com eles, recomenda-se, por fim, a leitura de um número razoável de passagens de texto que as contenham. Alguns trechos foram selecionados abaixo, em caráter ilustrativo:

\begin{quote}
\emph{"O que parece estarrecedor para quem nunca ouviu \textbf{Dilma} antes - e tenho colegas jornalistas que nunca a viram discursando ou dando \textbf{entrevista} - é absolutamente familiar para os frequentadores desta coluna. Que há nove meses têm acesso a veementes indícios, há muito transformados em provas documentais, de que \textbf{Dilma} é uma afronta imposta ao \textbf{Brasil} por \textbf{Lula}, num [sic] crime lesa-pátria sem perdão."}

{\small Retirado de "O som perturbador do neurônio em ebulição", da coluna de Augusto Nunes - 20/07/2010}
\end{quote}

\begin{quote}
\emph{"Que o \textbf{tucano} \textbf{José} Serra se saiu muito melhor no \textbf{Jornal} Nacional e que a eleição é, sim, de continuidade — no sentido de que não cabe mais falar em ruptura. E fiz uma crítica ou outra ao governo \textbf{Lula}."}

{\small Retirado de "A cabeça dos brasileiros... autoritários", do \emph{blog} de Reinaldo Azevedo - 15/08/2010}
\end{quote}

\begin{quote}

\emph{"A última bala na agulha do \textbf{Serra} é a baixaria. Só que, na era da internet, a baixaria – Lunus (para desconstruir Roseana Sarney) e aloprados do \textbf{PT} (para mandar as ambulâncias superfaturadas para o Inferno) – não tem o mesmo efeito do passado.É o caso dos aloprados do tal dossiê que ele vai ter que explicar na Justiça."}

{\small Retirado de "Serra só tem uma saída: pendurar FHC no pescoço", do \emph{site} Conversa Afiada - 07/06/2010}
\end{quote}

\begin{quote}
\emph{"A entrevista de \textbf{Dilma} ao JN foi didática: \textbf{Dilma} conseguiu colar sua \textbf{candidatura} como continuidade das políticas do governo \textbf{Lula}. Ponto pra ela.Por outro lado, o casal número um do JN da \textbf{Globo} escorregou e mostrou claramente contra quem trabalham em 2010 e a favor de quem se esforçam para mudar tudo o que está aí."}

{\small Retirado de "O povo não é (mais) bobo...", do \emph{blog} Luis Nassif Online - 10/08/2010}
\end{quote}



\chapter{Conclusão}
\label{conclusoes}

Com a disseminação da Web, a elaboração e disseminação de textos carregados com um ponto de vista se popularizou. Não se tratam de documentos que trazem opiniões pontuais a respeito de um único objeto, como um filme ou um livro, mas sim a exposição de argumentos e ideias que, unidos, transmitem a defesa de uma posição a respeito de um certo tema. A leitura da \emph{survey} de Pang e Lee e de alguns artigos, como o de Pang, Lee e Vaithyanathan sobre filmes no IMDb\footnote{http://imdb.com/} ou o de Dave, Lawrence e Pennock sobre produtos vendidos na Amazon\footnote{http://amazon.com/}, sugere que a classificação por opiniões tem sido mais explorada em \emph{datasets} que envolvem produtos ou serviços \cite{omsa} \cite{thumbs-up} \cite{peanut-gallery}. Os artigos revisados no Capítulo \ref{chap3}, por sua vez, indicam que a classificação por perspectivas é mais aplicada a \emph{datasets} que envolvem posicionamentos políticos e temáticas polêmicas, como pena de morte.

A revisão apresentada no Capítulo \ref{chap3}, além de indicar preferências temáticas na área, também aponta algumas predileções metodológicas. A maioria dos trabalhos classifica os documentos com um Naïve Bayes ou um SVM\footnote{Esses classificadores são apresentados no Capítulo \ref{basicos}.}. Em boa parte dos casos, essa classificação explora apenas as diferentes contagens de palavras nos documentos. Alguns outros artigos exploram propriedades sintáticas ou semânticas dos \emph{datasets} analisados. Outros, por fim, consideram as interações entre os documentos na determinação de seus pontos de vista. A contagem de palavras, simples de ser computada em qualquer língua, foi a característica mais estudada pelos trabalhos revisados - ainda que não tenha recebido destaque em todos eles. Em boa parte dos casos, seu uso exclusivo já é suficiente para uma boa classificação. Em alguns outros, o uso de outras estratégias, mais complexas, foi fundamental para a melhoria na classificação. Diante disso, recomenda-se, inicialmente, para o problema de classificação por perspectiva, o uso \emph{exclusivo} de contagens de palavras. Se os resultados obtidos não forem satisfatórios, recomenda-se a escolha de um subconjunto ótimo de palavras, como indica o trabalho de Durant e Smith \cite{durant-smith}. Apenas se não houver melhorias significativas, recomenda-se o uso de outras características, um pré-processamento de documentos mais refinado ou a investigação das interações entre os documentos. Essas recomendações, portanto, partem das ideias mais simples para as mais complexas. Dessa forma, evitam-se esforços desnecessários. 
%INTERNACIONALIZÁVEEEEEL!!
%o uso de um pré-processamento mais refinado dos documentos ou 

A contagem de palavras foi explorada por muitos artigos, como característica diferenciadora de perspectivas. De fato, a classificação baseada em contagens de palavras, ou em alguma de suas variações (como a presença/ausência de palavras), assume uma hipótese defendida em estudos de Linguística de corpus: a quantidade de vezes que uma palavra é mencionada em um documento (sua contagem) está diretamente relacionada a quais ideias ele destaca \cite{teubert}. Como consequência, ela funciona melhor em \emph{datasets} nos quais o emprego de palavras varia significativamente por perspectiva. A fim de investigar como o uso de palavras, refletido em suas contagens, se associa ao desempenho da classificação, o Capítulo \ref{chap3} apresenta alguns experimentos neste sentido. Dois corpora são classificados com um Naïve Bayes baseado em contagens de palavras. Via validação cruzada, as taxas de acerto obtidas foram, respectivamente, de 86.22\% e 54.01\%. Isso significa que, no primeiro corpus, as contagens de palavras foram suficientes para identificar corretamente a classe dos documentos de forma satisfatória; no segundo corpus, tem-se o caso contrário. Em outras palavras, isso significa que, no primeiro corpus, o emprego de palavras varia significativamente por perspectiva; no segundo, não. 

A fim de verificar quais palavras foram mais enfocadas por cada ponto de vista, ampliando a compreensão dos resultados obtidos com o Naïve Bayes, foram feitos experimentos com o modelo L-LDA\footnote{O L-LDA é apresentado no Capítulo \ref{basicos}.}. O modelo associa palavras a tópicos, facilitando a visualização de quais são mais relevantes para cada um deles. Nesses experimentos, cada perspectiva de cada \emph{dataset} corresponde a um tópico, e um último tópico, neutro, é associado a todos os documentos. A finalidade do tópico neutro é filtrar palavras muito comuns em cada corpus, independentemente de perspectiva. Com isso, evidencia-se quais palavras receberam mais destaque por cada perspectiva. Selecionando-se as dez palavras mais frequentemente associadas a cada perspectiva, já é possível visualizar que, no primeiro corpus, perspectivas diferentes enfatizam ideias mais distintas do que no segundo corpus. Isso amplia a compreensão das classificações feitas com o Naïve Bayes, apresentando cenários em que o uso exclusivo de contagens de palavras é ou não suficiente para uma boa classificação. É importante frisar que não se encontrou nenhum outro trabalho que faça uso de um L-LDA para compreender, ainda que parcialmente, como certos termos são enfocados por diferentes perspectivas.

Aproveitando as revisões e discussões exploradas no Capítulo \ref{chap3}, decidiu-se fazer um estudo de caso envolvendo um \emph{dataset} brasileiro. Não foi encontrado nenhum outro trabalho que aplique as técnicas dessa sub-área de Mineração de Opinião a documentos escritos em português - em particular, envolvendo conteúdo brasileiro. Considerando que 2010 é ano de eleições presidenciais no Brasil, a abundância de artigos que carregam pontos de vista típicos da oposição e da situação foi explorada. Foi construído um corpus sobre o atual governo e as eleições, composto de material coletado em colunas, \emph{sites} e \emph{blogs} mantidos por jornalistas de notoriedade nacional. Em seguida, esse corpus foi dividido entre as perspectivas pró-situação e pró-oposição, e foi classificado com um Naïve Bayes. As características dos documentos exploradas por esse classificador foram suas contagens de palavras. Os resultados obtidos com essa metodologia foram muito positivos: a taxa de acerto, em particular, foi de 89.43\%. Isso significa que esses jornalistas consolidam suas perspectivas sobre a política brasileira já nas palavras que destacam. De fato, alguns trechos elencados no Capítulo \ref{estudo}, no qual o estudo é apresentado, sugerem que seus pontos de vista são defendidos com muita veemência. 

Esse estudo de caso também investiga, assim como no Capítulo \ref{chap3}, quais palavras receberam mais destaque por cada perspectiva. A análise com um L-LDA, aplicado de forma semelhante àquela do Capítulo \ref{chap3}, indica que os jornalistas pró-situação dão mais ênfase ao candidato de oposição José Serra do que os pró-oposição. Esses últimos, por sua vez, dão maior destaque ao presidente Lula e à sua candidata, Dilma Rousseff - ambos correspondem à situação. Isso sugere que os jornalistas enfatizam o ataque aos candidatos que se opõem, ideologicamente, aos pontos de vista que eles defendem. Isso reforça a ideia de que a defesa de um posicionamento, muitas vezes, compreende o ataque a posicionamentos opostos, como sugerem Somasundaran e Wiebe em seu trabalho sobre debates \emph{online} \cite{wiebe}.  

\section{Dificuldades encontradas}

A classificação de documentos de acordo com seus pontos de vista sobre um tema é um problema relativamente novo. De fato, ele só foi estabelecido como sub-área da Mineração de Opinião em 2008, na \emph{survey} de Pang e Lee. Por este motivo, para entender melhor o problema, foi necessário buscar artigos em diversas conferências que envolviam a área de Mineração de Opinião. A filtragem de quais resultados realmente tinham a ver com o tema dessa monografia não foi exatamente difícil, mas envolveu um trabalho manual considerável. Para o estudo de caso, a extração, filtragem e padronização dos documentos que compõem o corpus também envolveu algum trabalho manual. Embora essas tarefas tenham sido realizadas de forma automatizada, a construção dos \emph{scripts} dependia da compreensão de como o conteúdo estava disposto em cada veículo.  

\section{Trabalhos futuros}

Futuramente, pretende-se estender o estudo de caso do Capítulo \ref{estudo} a textos políticos escritos por cidadãos comuns em seus \emph{blogs}, o que pode contribuir positivamente para a compreensão de como o brasileiro se posiciona politicamente na Web. Além disso, o estudo também deve ser ampliado para identificar as perspectivas contidas nos comentários feitos aos artigos do corpus, em seus \emph{blogs}, \emph{sites} e colunas, a fim de se avaliar como eles refletem o posicionamento dos leitores em relação àquilo que leram. Este tipo de análise pode ajudar a compreender o impacto destes artigos em seus leitores e a formação de posicionamentos na mídia brasileira \emph{online}. Caso essas tarefas de classificação não sejam bem resolvidas com o uso exclusivo de contagens de palavras, esforços direcionados na busca de características mais complexas, envolvendo aspectos sintáticos/semânticos dos documentos, devem ser feitos. É válido ressaltar que isso provavelmente implicaria no desenvolvimento de ferramentas gramaticais de apoio, voltadas para a língua portuguesa.




\bibliography{monografia}
\appendix

\chapter{Teorema de Bayes e notações}

A Estatística se alia à Teoria da Probabilidade para estimar e analisar a ocorrência de eventos, como \emph{Ganho de cem reais em sorteio} ou \emph{Chuva ao meio-dia}. Dados dois eventos \emph{A} e \emph{B}, temos que a probabilidade \emph{a priori} de \emph{A} acontecer \textbf{ignora} a ocorrência do evento \emph{B}. Essa probabilidade é representada pela notação \ensuremath{P(A)} \cite{spiegelhalter}. Analogamente, para o evento \emph{B}, tem-se \ensuremath{P(B)}. Na prática sabe-se, intuitivamente, que a ocorrência de determinados eventos interfere no acontecimento de outros. Se às onze e meia da manhã observa-se o evento \emph{Céu nublado}, a probabilidade do evento \emph{Chuva ao meio-dia} ocorrer pode diferir daquela que ignora esse primeiro evento. Neste sentido, a probabilidade de um evento \emph{A} ocorrer \emph{dado} que \emph{B} ocorreu recebe o nome de probabilidade condicional de \emph{A} dado \emph{B}, e é denotada por \ensuremath{P(A\mbox{ }|\mbox{ }B)} \cite{spiegelhalter}. Analogamente, tem-se \ensuremath{P(B\mbox{ }|\mbox{ }A)}. O Teorema de Bayes correlaciona essas probabilidades da seguinte forma \cite{spiegelhalter}

\begin{equation}
\label{bayes-apendice}
\ensuremath{P(A\mbox{ }|\mbox{ }B) = \frac{P(B\mbox{ }|\mbox{ }A)P(A)}{P(B)}} 
\end{equation}

No contexto da Equação \ref{bayes-apendice}, \ensuremath{P(A\mbox{ }|\mbox{ }B)} recebe o nome de probabilidade \emph{a posteriori} de \emph{A} e \ensuremath{P(B\mbox{ }|\mbox{ }A)} recebe o nome de \emph{likelihood} \cite{spiegelhalter}. O classificador Naïve Bayes se baseia em uma aplicação direta do Teorema de Bayes.  Dados os eventos \emph{Obtenção de um documento x} e \emph{Obtenção de uma classe y}, o classificador deve estimar a ocorrência do segundo evento assumindo que o primeiro já ocorreu.   

\chapter{Código do Naïve Bayes}

A implementação do Naïve Bayes foi feita em Python, aplicando o algoritmo de inferência Gibbs Sampling para estimar parâmetros e classes. Há uma classe responsável pelo \emph{parsing} dos documentos, \textbf{CorpusParser}, e outra para o classificador em si, \textbf{NaiveBayesSampler}.

\section{CorpusParser}

\begin{lstlisting}

#Aline Bessa - 17/05/2010
import sys
import os
import numpy as np
import random
from nltk import *

stop_words = [] 

def read_file(file):
   text = []
   line = file.readline()
   while(line != ""):
     text.append(str(line).strip())
     line = file.readline()
   return text

def get_words(text):
  import re
  wre = re.compile(r"(\w)+")
  l = 0
  while l < len(text):
       s = wre.search(text,l)
       try:
           yield text[s.start():s.end()]
           l = s.end()
       except:
           break

class CorpusParser(object):
  def __init__(self, train_corpus, test_corpus):
    self.train_corpus = train_corpus
    self.Ntrain_corpus = len(os.listdir(train_corpus))
    self.test_corpus = test_corpus
    self.Ntest_corpus = len(os.listdir(test_corpus))
    self.train_list_docs = os.listdir(train_corpus)
    self.test_list_docs = os.listdir(test_corpus)
    self.train_docs = []
    self.test_docs = []

    self.all_words = []

  def train_dirname(self):
    return self.train_corpus
  def test_dirname(self):
    return self.test_corpus

  def chop(self, raw_sntc):
    x = raw_sntc.find(">") + 1
    y = raw_sntc.find("</")
    return raw_sntc[x:y]

  def words(self):
    return self.all_words

  def get_new_words(self, doc):
    for i in xrange(len(doc)):
      for t in get_words(str(doc[i])):
        t = t.lower().strip()
       # t = en.noun.singular(t)
        if not t in self.all_words and not t in stop_words:
          self.all_words.append(t)

  def ldocs(self):
    return self.list_docs

  def dirname(self):
    return self.corpus

  def get_doc(self, doc):
    counts = np.zeros(len(self.all_words))
    for sntc in doc:
      for t in get_words(sntc):
        t = t.lower().strip()
       # t = en.noun.singular(t)
        if not t in stop_words and t in self.all_words:
          counts[self.all_words.index(t)] += 1.
    return counts

  #TODO cut redundancy between docs and dir params
  def iterate_corpus(self, docs, dir, action, 
                     doc_preprocessing = None):
    for file in docs:
      f = open(os.path.realpath(dir + '/' + file), 'r')
      doc = read_file(f)
      if doc_preprocessing:
        action(doc_preprocessing(doc))
      else:
        action(doc)
      f.close()


  ###get every single word first
  def pdocs(self):

    self.iterate_corpus(self.train_list_docs, self.train_corpus, 
                        self.get_new_words)
    self.iterate_corpus(self.test_list_docs, self.test_corpus, 
                        self.get_new_words)
    self.iterate_corpus(self.train_list_docs, self.train_corpus, 
                        self.train_docs.append, self.get_doc)
    self.iterate_corpus(self.test_list_docs, self.test_corpus, 
                        self.test_docs.append, self.get_doc)
    return self.train_docs, self.test_docs
 # print self.all_words, len(self.all_words)

if __name__ =='__main__':
  a = CorpusParser(sys.argv[1], sys.argv[2])
  a.pdocs()


\end{lstlisting}

\section{NaiveBayesSampler}

\begin{lstlisting}

#Aline Bessa - 27/04/2010
#Naive Bayes modelling for sentiment analysis in docs

import numpy as np
import random
import math
from parse_nb import CorpusParser
import sys
import os

class NaiveBayesSampler(object):
  def __init__(self, a):
    pdocs = a.pdocs()
    self.train_docs = pdocs[0]
    self.test_docs = pdocs[1]
    self.docs = self.train_docs + self.test_docs

    if not self.docs: 
      self.all_words = 0
    else:
      self.all_words = len(self.docs[0])
    #print self.all_words
    self.train_dirname = a.train_dirname()
    self.test_dirname = a.test_dirname()
    self.labels = [] 
    self.Ntraindocs = len(self.train_docs)
    self.Ntestdocs = len(self.test_docs)

  def initial_label(self):
    return np.random.binomial(1, self.pi)

  def get_real_label(self, dir, index):
    #assumptions, assumptions: sit = 1; opos = 0
    ldocs = os.listdir(dir)
    if ldocs[index].find("opos") >= 0:
      return 0
    else:
      return 1

  def label_documents(self):
    for i in xrange(self.Ntraindocs):
      label = self.get_real_label(self.train_dirname, i)
      self.labels.append(label)
    for i in xrange(self.Ntestdocs):
      label = self.initial_label()
      self.labels.append(label)

  def pLi(self, j, doc):
    #This is how the probabilities are 
    #ACTUALLY calculated (Gibbs Sampling)   
    # Log - implementation detail: it smoothes the computations
    t1 = np.log((self.dccs[j] + self.Gammapi[j])/
        (len(self.labels) + self.Gammapi[1] + self.Gammapi[0] -1))
    t2 = doc*np.log(self.theta[j])
    t2 = np.sum(t2)
    return t1+t2


  def pick_label(self, j):
    lclass = self.labels[j]
    self.dccs[lclass] -= 1 
    self.wccs[lclass] -= self.docs[j] 

    pL0 = self.pLi(0, self.docs[j]) 
    pL1 = self.pLi(1, self.docs[j])
    loglr = pL1-pL0
    lr = np.exp(loglr)
    if lr == np.inf:
      label = 1
    else:
      if np.random.binomial(1, lr/(1+lr)):
        label = 1
      else:
        label = 0
    self.labels[j] = label
    self.dccs[label] += 1
    self.wccs[label] += self.docs[j]
    return label
  
  def get_denominator(self):
      den = np.zeros(2)
      for i in range(self.Ntestdocs):
        den[self.get_real_label(self.test_dirname, i)] += 1
      return den

  def num_den(self):
    den = self.get_denominator()
    num = np.zeros(2)
    
    for i in range(self.Ntestdocs):
      #shifting Ntraindocs positions
      l = self.labels[i + self.Ntraindocs] 
      if l == self.get_real_label(self.test_dirname, i):
        num[self.labels[i + self.Ntraindocs]] += 1
    
    return num, den 

  def precision(self, num, den):
    #TP/(TP+FP)
    return num[0]/(num[0] + (den[1] - num[1])), 
           num[1]/(num[1] + (den[0] - num[0]))

  def accuracy(self, num, den):
    return num[0]/den[0], num[1]/den[1], 
          (num[0] + num[1])/(den[0] + den[1])
 
  def sample(self, nsamples):
    self.Gammapi = np.array([1., 1.])
    self.Gammatheta = np.array([1. for i in xrange(self.all_words)])
    self.pi = random.betavariate(1,1)
    self.theta = np.zeros((2, self.all_words)) 
    self.theta[0] = np.random.dirichlet(self.Gammatheta) 
    self.theta[1] = np.random.dirichlet(self.Gammatheta) 
    self.label_documents()
    self.wccs = np.zeros((2, self.all_words))
    self.dccs = np.zeros(2)

    for i, a in enumerate(self.labels):
      self.dccs[a] += 1
      self.wccs[a] += self.docs[i]

    for i in xrange(nsamples):
      for j in xrange(self.Ntestdocs):
        #shifting Ntraindocs positions
        self.pick_label(j + self.Ntraindocs) 
      self.theta[0] =
      np.random.dirichlet(self.Gammatheta + self.wccs[0])
      self.theta[1] = 
      np.random.dirichlet(self.Gammatheta + self.wccs[1])

      num, den = self.num_den()
      print num[0], num[1]
      print den[0], den[1]
      print "accuracy"
      ac = self.accuracy(num, den)
      print ac
      print "precision"
      pr = self.precision(num, den)
      print pr
      print "f1 measure"
      print 2*pr[0]*ac[0]/(pr[0] + ac[0]), 
            2*pr[1]*ac[1]/(pr[1] + ac[1]) 

if __name__=='__main__':
  a = CorpusParser(sys.argv[1], sys.argv[2]) 
  s = NaiveBayesSampler(a)
  s.sample(500)

\end{lstlisting}


\end{document}

